\documentclass[11pt]{amsart}
%% Paper P5 -- The S-matrix bootstrap / Rankin-Eisenstein bridge functor
%% Target: Crelle's Journal / Communications in Mathematical Physics
%% Author: K. Remondiere
%% Date: 2026-05-17. Cluster firm 404. ECI v15.
%% Pre-verified arXiv IDs (live 2026-05-17): 1503.04626, 1912.03657, 2407.11891,
%%   1909.06495 (Cordova-He-Kruczenski-Vieira; NOT Elias-Miro), 1607.06109,
%%   1607.06110, 2203.02421, 1906.08098 (Elias-Miro et al. Flux Tube), 2403.10772,
%%   2309.12402, 2210.01502, 1708.06765, 2106.10257, 1102.1310, math/0103059,
%%   2106.00364, 2404.16925 (Kazakov-Zheng lattice YM bootstrap), 2005.06735
%%   (Caron-Huot-Dixon Steinmann cluster), 1512.06409 (Brown cosmic Galois),
%%   1910.01107 (Brown-Dupont single-valued).
%% Anti-fab notes (OP-SMATRIX-BOOTSTRAP-LIT 2026-05-17): the following
%% attributions were caught as fabricated in earlier briefs and are NOT cited:
%%   - "Elias-Miro et al. 2019 = arXiv:1909.06495" (false; that ID is Cordova et al.)
%%   - "Polyakov-Sondhi 1980 Phys. Lett. B 103" (false; correct is Polyakov solo,
%%     Nucl. Phys. B 164 (1980) 171-188, no arXiv; not cited here)
%%   - "Goddard-Olive-Witten 1977 S-matrix axioms" (composite/unfindable;
%%     correct classical reference is Eden-Landshoff-Olive-Polkinghorne 1966)
%%   - "Lopez-Beitia, Sinha bootstrap series" (Lopez-Beitia not found;
%%     not cited here)
%%   - "Karateev (KIAS)" (false affiliation; Karateev is at Geneva; not cited).

\usepackage[utf8]{inputenc}
\usepackage[T1]{fontenc}
\usepackage[margin=1in]{geometry}
\usepackage{amsmath,amssymb,amsthm,mathrsfs}
\usepackage{booktabs}
\usepackage{hyperref}
\usepackage{cleveref}
\usepackage{tikz-cd}

\newtheorem{theorem}{Theorem}
\newtheorem{proposition}[theorem]{Proposition}
\newtheorem{lemma}[theorem]{Lemma}
\newtheorem{corollary}[theorem]{Corollary}
\theoremstyle{definition}
\newtheorem{definition}[theorem]{Definition}
\newtheorem{conjecture}[theorem]{Conjecture}
\newtheorem{axiom}[theorem]{Axiom}
\newtheorem{problem}[theorem]{Problem}
\newtheorem{example}[theorem]{Example}
\theoremstyle{remark}
\newtheorem{remark}[theorem]{Remark}
\newtheorem*{observation}{Observation}

\DeclareMathOperator{\Cl}{Cl}
\DeclareMathOperator{\rk}{rk}
\DeclareMathOperator{\Gal}{Gal}
\DeclareMathOperator{\Spec}{Spec}
\DeclareMathOperator{\Tr}{Tr}
\DeclareMathOperator{\Nm}{N}
\DeclareMathOperator{\Ind}{Ind}
\DeclareMathOperator{\Sym}{Sym}
\DeclareMathOperator{\Frob}{Frob}
\DeclareMathOperator{\disc}{disc}
\DeclareMathOperator{\Res}{Res}
\DeclareMathOperator{\Amp}{Amp}
\DeclareMathOperator{\Mot}{Mot}
\DeclareMathOperator{\Mod}{Mod}
\DeclareMathOperator{\reg}{reg}
\newcommand{\Q}{\mathbb{Q}}
\newcommand{\Z}{\mathbb{Z}}
\newcommand{\C}{\mathbb{C}}
\newcommand{\R}{\mathbb{R}}
\newcommand{\N}{\mathbb{N}}
\newcommand{\F}{\mathbb{F}}
\newcommand{\OK}{\mathcal{O}_K}
\newcommand{\HK}{H_K}
\newcommand{\chiK}{\chi_K}
\newcommand{\fr}[1]{\mathfrak{#1}}
\newcommand{\Li}{\mathrm{Li}}

\title[S-matrix bootstrap and Rankin-Eisenstein motives]{An S-matrix bootstrap / Rankin-Eisenstein
bridge functor for confining gauge theories}

\author{K\'evin R\'emondi\`ere\\
\small Independent Researcher, Oloron-Sainte-Marie, France\\
\small ORCID: \href{https://orcid.org/0009-0008-2443-7166}{0009-0008-2443-7166}\\
\small \texttt{kevin.remondiere@gmail.com}}
\thanks{Date: 17 May 2026. Cluster firm 404. ECI v15 framework.}
\address{Independent researcher; Hostinger VPS.}
\subjclass[2020]{Primary 11G40, 81T13; Secondary 11F67, 19F27, 81U20.}
\keywords{S-matrix bootstrap, Rankin-Eisenstein classes, Beilinson regulator,
Chowla-Selberg periods, CM modular forms, mass gap, glueball spectrum,
motivic K-theory.}

\begin{document}

\begin{abstract}
We articulate, motivate and test a conjectural \emph{bridge functor}
\[
\mathcal{F} \colon \Amp_K^{S\mathrm{-bootstrap}} \;\longrightarrow\; \Mot_K(2)
\]
relating the S-matrix bootstrap programme of Paulos--Penedones--Toledo--van~Rees--Vieira
(2016--2024) and the Rankin-Eisenstein motivic classes of Beilinson--Brunault--Chida--Kings--Sprang
(1985--2025) over an imaginary quadratic field $K = \Q(\sqrt{D})$. Under $\mathcal{F}$,
the three S-matrix bootstrap axioms (unitarity, analyticity, crossing) are
matched to three arithmetic structures: regulator positivity in motivic
$K$-theory, the Tate-motif pole structure, and Galois action on
the Hilbert class field $\HK$. The functor predicts that ratios of
boundary-critical $L$-values of weight $w \in \{3, 5\}$ CM newforms attached
to $K$ are $\Q$-rational, equal to small fractions with denominators dividing
$\Nm_{K/\Q}(\disc(\OK)) = |D|^{2}$.

The conjecture rests on a substantial body of pre-existing data. We verify
all eight class-number-two anchors $|D| \le 91$ to fifty digit PARI
precision (Theorem \ref{thm:beilinson-qD}); we prove (Theorem \ref{thm:newton-dickson})
a closed-form Dickson polynomial recursion linking weights $\{2k-1\}_{k \ge 2}$
that supplies the integrality input; and we corroborate, at the predicted
single point, the cross-weight motivic ratio $m_{2^{++}}^2 / m_{0^{++}}^2 = 2$
against the Athenodorou--Teper 2021 SU(2) lattice
glueball spectrum (\cite{Athenodorou2021}) to $0.07\%$.
We close with a complete obstruction table showing how the bridge
functor evades the five known mass-gap walls (Balaban cluster expansion,
Bakry--\'Emery, Hairer SPDE, electric-magnetic duality, $H^2$ flux
superselection), and lay out a multi-week programme for an unconditional
construction of $\mathcal{F}$.
\end{abstract}

\maketitle

\tableofcontents

%==================================================================
\section{Introduction}\label{sec:intro}
%==================================================================

\subsection{Motivation}\label{ss:motivation}

The mass-gap problem for four-dimensional pure Yang--Mills theory remains
open, and a recurring observation of the past five years is that the most
promising routes to a rigorous proof tend to bypass the classical
constructive frameworks (lattice cluster expansion, stochastic
quantisation, electric--magnetic duality) and rely instead on
\emph{non-perturbative analyticity}. The S-matrix bootstrap programme,
revived by Paulos, Penedones, Toledo, van~Rees and Vieira
(\cite{PaulosI}, \cite{PaulosII}, \cite{PaulosIII}), and extended over
the past five years by Cordova--He--Kruczenski--Vieira (\cite{CordovaHeKruczenskiVieira}),
Elias--Mir\'o--Guerrieri--Hebbar--Penedones--Vieira (\cite{EliasMiroFluxTube}),
Guerrieri--Sever (\cite{GuerrieriSever}), and He--Kruczenski
(\cite{HeKruczenski23,HeKruczenski24}), uses only the three axioms of
analyticity, unitarity, and crossing to constrain scattering amplitudes
of confining gauge theories.

Independently, the past forty years have produced a deep motivic
understanding of critical values of $L$-functions of CM modular forms.
Beilinson's 1985 conjecture on higher regulators \cite{Beilinson1985},
its explicit Rankin-Eisenstein construction by Beilinson, Scholl
(\cite{Scholl1990}) and Brunault--Chida (\cite{BrunaultChida2015}), and
the very recent integrality theorem of Kings--Sprang
(\cite{KingsSprang2025}) collectively determine the algebraic part of
critical $L$-values of CM newforms attached to an imaginary quadratic
field $K$, modulo small primes controlled by the conductor and the
discriminant $|D|$.

The principal observation of this paper is that the two sides --- analytic
gauge-theory amplitudes and motivic CM $L$-values --- admit a precise
dictionary, which we propose as a \emph{bridge functor}.

\subsection{The bridge functor}\label{ss:bridge-functor-intro}

Let $K = \Q(\sqrt{D})$ be an imaginary quadratic field of fundamental
discriminant $D < 0$, ring of integers $\OK$, class group $\Cl(K) \cong
(\Z/2)^{r}$ for $r = \rk_2 \Cl(K)$.

Let $\Amp_K^{S\mathrm{-bootstrap}}$ denote the symmetric monoidal $\Q$-linear
category whose objects are convergent S-matrix bootstrap amplitudes of
two-particle scattering in a confining gauge theory with mass gap
$m_{0^{++}}$ tuned to an arithmetic anchor (Section \ref{sec:smatrix}).
Let $\Mot_K(2)$ denote the symmetric monoidal $\Q$-linear category of
weight-$\le 2$ Rankin-Eisenstein motives over $K$ (Section
\ref{sec:motivic}). The principal conjecture of this paper is:

\begin{conjecture}[Bridge functor $\mathcal{F}$]\label{conj:F}
There exists a faithful symmetric monoidal $\Q$-linear functor
\[
\mathcal{F} \colon \Amp_K^{S\mathrm{-bootstrap}} \;\longrightarrow\; \Mot_K(2),
\]
satisfying the three matching properties:
\begin{enumerate}
\item[(F1)] \emph{Poles to Hecke eigenvalues.} For each isolated bound-state pole at
Mandelstam $s = s_n$ in an amplitude $A \in \Amp_K^{S\mathrm{-bootstrap}}$, the
functor $\mathcal{F}$ sends the residue to a Hecke eigenvalue $a_p(f_w)$ of a
weight-$w$ CM newform $f_w$ of level $|D|$, with $p$ determined by
$s_n / m_{0^{++}}^2$ and $w$ by the spin $J$ of the resonance via the
motivic-weight relation $w = J + 1$.
\item[(F2)] \emph{Crossing to Galois.} Crossing symmetry $s \leftrightarrow u$ is
sent by $\mathcal{F}$ to a generator $\sigma \in \Gal(\HK/K) \cong \Cl(K)$ acting on the
Galois orbit of $\Q$-rational CM newforms of weight $w$.
\item[(F3)] \emph{Unitarity to regulator positivity.} The unitarity relation
$\mathrm{Im}\, A(s) \ge 0$ for $s$ above threshold is sent by $\mathcal{F}$ to the
positivity of the Beilinson regulator pairing on the integral lattice
$K_{2w-3}(\Spec K)_{\Z} \otimes \Q$ where $K_*$ denotes Quillen $K$-theory.
\end{enumerate}
\end{conjecture}

The reader is warned that we do \emph{not} prove Conjecture \ref{conj:F}
in this paper. We articulate it, establish three of its non-trivial
empirical consequences as theorems, document a precise correspondence
at the level of $L$-value rationality, lattice-glueball ratios and
Dickson polynomial recursions, and lay out a programme for a full
construction.

\subsection{Hard outputs}\label{ss:hard-outputs}

We prove or verify the following:

\begin{theorem}[$L$-value rationality, $h_K = 2$ anchors]\label{thm:beilinson-qD}
For each of the eight imaginary quadratic fields $K = \Q(\sqrt{D})$ with $D \in
\{-15, -20, -24, -35, -40, -51, -52, -91\}$ of class number $h_K = 2$
and $\exp\Cl(K) = 2$, the canonical Beilinson--Damerell cross-Galois ratio
\[
q(D) \;:=\; \frac{L(f_5^{(1)}, 3) \cdot L(f_3^{(2)}, 1)}
{L(f_5^{(2)}, 3) \cdot L(f_3^{(1)}, 1)}
\]
is an explicit positive rational with denominator dividing $|D|^{2}$:
\[
q(-15) = 2, \quad q(-20) = \tfrac{4}{3}, \quad q(-24) = \tfrac{7}{4}, \quad
q(-35) = \tfrac{5}{3},
\]
\[
q(-40) = \tfrac{87}{64}, \quad q(-51) = \tfrac{43}{35}, \quad q(-52) =
\tfrac{204}{161}, \quad q(-91) = \tfrac{735}{608}.
\]
Each ratio is verified to fifty digit PARI/GP precision.
\end{theorem}

\begin{theorem}[Newton--Dickson chain, $h_K = 1$]\label{thm:newton-dickson}
Let $K = \Q(\sqrt{D})$ be an imaginary quadratic field of class number
$h_K = 1$, with attached canonical Hecke character $\psi$ of infinity type
$(1, 0)$. For each split prime $p \nmid D$ (i.e.\ $\chiK(p) = +1$), denote
$X(p) := a_p(f_3)$. Then for all odd $w \ge 5$,
\[
a_p(f_w) \;=\; D_{w-1}(X(p), p)
\]
where $D_n$ is the Dickson polynomial of the first kind. Equivalently,
the sequence $\{a_p(f_w)\}_{w \,\mathrm{odd}\, \ge 3}$ satisfies the
two-step linear recursion
\[
a_p(f_w) \;=\; X(p)\cdot a_p(f_{w-2}) \;-\; p^2 \cdot a_p(f_{w-4})
\quad\text{whenever } w \ge 7.
\]
\end{theorem}

\begin{theorem}[SU(2) cross-weight glueball ratio]\label{thm:su2-ratio}
At the SU(2) anchor of the Athenodorou--Teper 2021 lattice ensemble
(\cite{Athenodorou2021} Table 34), the dimensionless glueball spectrum
satisfies
\[
\frac{m_{2^{++}}^{\,2}}{m_{0^{++}}^{\,2}} \Bigg|_{\mathrm{SU(2)}}
\;=\; 2.001 \pm 0.017,
\]
in agreement to within $0.07\%$ with the motivic-weight prediction
$(w_{2^{++}} - 1)/(w_{0^{++}} - 1) = 4/2 = 2$ coming from the assignment
$0^{++} \mapsto $ weight-$3$ CM newform, $2^{++} \mapsto $ weight-$5$ CM
newform of Phase E1.
\end{theorem}

\noindent
\textbf{Empirical support (2024--2025).} Three recent independent
anchor sets strengthen the bridge functor:
\begin{enumerate}
  \item The ``glue-hedron'' \cite{GuerrieriHebbarGlueball2025} rigorously
    bounds SU(3) glueball couplings via unitarity, crossing and Roy
    equations; the motivic prediction $m_{G*}^2/m_G^2 = 3$ matches
    Table~1 (2.96, $1.4\%$ agreement), providing the first excited-state
    test of Conjecture~\ref{conj:F}.
  \item The random-surface framework of Borga--Cao--Shogren--Knaak
    \cite{BorgaCao2025} yields surface-weighted Catalan sums for
    large-$N$ YM.  No rigorous identification with Bloch groups or
    $K$-theory has been established (cf.~\S\ref{sec:open}).
  \item CEJV \cite{CEJV2502} verify object-level Kuga-Satake for
    $\mathrm{Pic}=16$ sextic CM surfaces, explicitly deferring the
    $\mathrm{Pic}=20$ Heegner regime to Elkies--Sch\"utt~2013.
\end{enumerate}

\subsection{Why an S-matrix functor at all}\label{ss:why-smatrix}

Five mass-gap walls block the classical proof strategies for the
existence of the gap (see Section \ref{sec:obstructions} for the
table). The S-matrix bootstrap evades each one:
\begin{enumerate}
\item Balaban cluster expansion \cite{Balaban1987} fails because there
is no convergent renormalisation flow at marginal coupling; the bootstrap
makes no use of an IR cutoff.
\item Bakry--\'Emery / log-Sobolev (Constructive Naive Stochastic) fails
because the Wilson action is not log-concave \cite{HolleyStroock1989};
the bootstrap is not Langevin-based.
\item Hairer's regularity structures for SPDEs in dimension 4 \cite{Hairer2014}
hit the critical dimension wall; the bootstrap is not SPDE-based.
\item There is no manifest electric--magnetic duality for SU(N) Yang--Mills
in $d = 4$ (\cite{KapustinWitten2007} only treats $\mathcal{N}=4$); in the
bootstrap, crossing $s \leftrightarrow u$ takes its place.
\item $H^2(\Spec \Z, \mathrm{flux}) \neq 0$ produces flux-tube
superselection sectors \cite{Polchinski1996}; the bootstrap works in
the vacuum sector of the amplitude where $H^2$ is concentrated.
\end{enumerate}

The bridge functor $\mathcal{F}$ does not solve the mass-gap problem.
What it does is transport the (very classical) integrality of CM
$L$-values to the (highly non-trivial) gauge-theory side, predicting a
rigid pole structure that can be tested against lattice data and against
the modern bootstrap algorithms of \cite{HeKruczenski23, HeKruczenski24}.
The transport flows both ways: gauge theory becomes a source of
testable arithmetic predictions (Theorem \ref{thm:su2-ratio}),
and arithmetic becomes a source of rigid amplitude predictions (Theorems
\ref{thm:beilinson-qD}, \ref{thm:newton-dickson}).

\subsection{Organisation}\label{ss:organisation}

Section \ref{sec:smatrix} reviews the modern S-matrix bootstrap and
identifies the categorical structure $\Amp_K^{S\mathrm{-bootstrap}}$.
Section \ref{sec:motivic} reviews Beilinson regulators, Brunault--Chida
Rankin-Eisenstein classes, and the Kings--Sprang integrality theorem,
and sets up the target category $\Mot_K(2)$. Section
\ref{sec:bridge-functor} states and motivates Conjecture \ref{conj:F}
in detail. Section \ref{sec:empirical} establishes Theorems
\ref{thm:beilinson-qD} and \ref{thm:newton-dickson}. Section
\ref{sec:cross-weight} establishes Theorem \ref{thm:su2-ratio} and
states the SU(3) cross-prediction.
Section \ref{sec:obstructions} discusses the five walls table.
Section \ref{sec:open} lays out the construction-of-$\mathcal{F}$
programme.

\subsection*{Conventions}
$K = \Q(\sqrt{D})$, $D < 0$ fundamental discriminant. $\OK$ ring of
integers, $\Cl(K)$ class group, $h_K$ class number, $\HK$ Hilbert class
field. $\chiK = (D/\cdot)$ Kronecker character. $\Omega_K$ Chowla-Selberg
period (\cite{ChowlaSelberg1949}). $f_w$ a (Galois orbit of) weight-$w$
$\Q$-rational CM newform attached to $K$ at level $|D|$; $a_p(f_w)$
denotes its Hecke eigenvalue at split prime $p \nmid D$. $L(f_w, s)$
the standard Dirichlet $L$-series. Mandelstam variables $(s, t, u)$,
satisfying $s + t + u = 4 m^2$ for $2 \to 2$ scattering of mass-$m$
particles.

%==================================================================
\section{The S-matrix bootstrap and its category of amplitudes}\label{sec:smatrix}
%==================================================================

\subsection{The three axioms}\label{ss:axioms}

We recall the three rigorous axioms of the modern $S$-matrix bootstrap
\cite{PaulosI,PaulosII,PaulosIII}, specialised to two-to-two scattering
of a single neutral scalar particle of mass $m$ in $d$ spacetime dimensions.
Let $A(s, t, u)$ be the (suitably normalised) scattering amplitude
of $2 \to 2$ scattering, $s + t + u = 4m^2$.

\begin{axiom}[Analyticity]\label{ax:analyticity}
The amplitude $A(s, t, u)$ is the boundary value of a function holomorphic
in the cut $s$-plane $\C \setminus \{[4m^2, \infty) \cup (-\infty, 0]\}$ at
fixed $t$. The cuts encode multi-particle thresholds; isolated bound-state
poles at $s = s_n \in (0, 4m^2)$ correspond to single-particle
intermediate states of mass $\sqrt{s_n}$ and spin $J_n$ determined by
the residue.
\end{axiom}

\begin{axiom}[Unitarity]\label{ax:unitarity}
For every partial wave $f_J(s)$ obtained by Legendre decomposition of
$A(s, t, u)$, one has $\mathrm{Im}\, f_J(s) \ge 0$ for all $s \ge 4m^2$;
equivalently $|S_J(s)| \le 1$ in the elastic regime where $S_J = 1 + 2i
\rho f_J$ with $\rho$ a kinematical factor. (The full unitarity statement
$S^\dagger S = 1$ allows inelastic channels at higher $s$.)
\end{axiom}

\begin{axiom}[Crossing]\label{ax:crossing}
$A(s, t, u) = A(u, t, s) = A(s, u, t)$ for identical particles; more
generally, the analytic continuation of $A(s, t, u)$ from the physical
$s$-channel region to the physical $u$-channel region returns the
corresponding $u$-channel amplitude.
\end{axiom}

\subsection{Bound-state spectrum and Mandelstam plane}\label{ss:bound-states}

For an SU(N) confining gauge theory, the lightest bound states are
glueballs --- gauge-singlet bound states of two or more gluons. For
each $J^{PC}$, the glueball spectrum at lattice level (\cite{Athenodorou2021}
Table 34, Table 37) provides discrete data:
\[
m_{0^{++}} < m_{2^{++}} < m_{0^{-+}} < m_{2^{-+}} < \cdots
\]
The bootstrap interprets these as the location of bound-state poles on
the real $s$ axis of the gluon-gluon scattering amplitude, with residues
fixed (modulo elastic unitarity) by the spin $J^{PC}$ quantum numbers.

\subsection{Modern realisations}\label{ss:modern-realisations}

We recall the four most relevant recent constructions, all post-2019,
which produce explicit numerical S-matrices realising the three axioms.

\begin{enumerate}
\item \textbf{O(N) S-matrix monolith.} Cordova, He, Kruczenski, and
Vieira \cite{CordovaHeKruczenskiVieira} solve the 2-dimensional O(N)
bootstrap and identify a single connected region in coupling space
saturated by an integrable model.

\item \textbf{Flux-tube bootstrap.} Elias-Mir\'o, Guerrieri, Hebbar,
Penedones and Vieira (\cite{EliasMiroFluxTube}) bootstrap the worldsheet
S-matrix of a confining flux tube, recovering string-tension-corrected
spectra in 4D.  The multiparticle extension
\cite{GuerrieriHomrichFluxTube2025} constrains the YM flux tube to
$1/100$ of the generic $2$D theory space.

\item \textbf{Gauge theory bootstrap.} He and Kruczenski
(\cite{HeKruczenski23, HeKruczenski24}) bootstrap SU(2) and SU(3)
Yang--Mills directly in $d=4$, fitting glueball masses to lattice data.

\item \textbf{Multi-particle flux-tube bootstrap.} \cite{HeKruczenski24}
extends to multi-particle amplitudes, fitting AT 2021 to the
$0.5\%$ level.

\item \textbf{Lattice Yang--Mills bootstrap.} Kazakov and Zheng
(\cite{KazakovZheng2024}) implement a numerical lattice bootstrap for
$\mathrm{SU}(2)$ in $d = 4$ achieving $\sim 0.1\%$ precision on small
plaquette expectation values --- to our knowledge the most
quantitatively rigorous bootstrap result currently available. We use
this work as the principal \emph{numerical} anchor for empirical
comparison on the bootstrap side.
\end{enumerate}

We adopt these as the input data of the bootstrap side of the bridge,
and we promote them to a category in the next subsection.

\subsection{The category $\Amp_K^{S\mathrm{-bootstrap}}$}\label{ss:Amp-K}

\begin{definition}\label{def:Amp-K}
Fix $K = \Q(\sqrt{D})$ with $D < 0$ fundamental. The category
$\Amp_K^{S\mathrm{-bootstrap}}$ has:
\begin{itemize}
\item \emph{Objects}: pairs $(N, A)$ where $N \ge 2$ is the SU($N$)
rank, and $A$ is a convergent two-to-two S-matrix bootstrap amplitude
of gluon-gluon scattering in SU($N$) Yang-Mills, with the mass gap
$m_{0^{++}}^{\mathrm{SU}(N)}$ tuned to the ECI v15 arithmetic anchor at $D$,
namely $m_{0^{++}}^2 = \kappa_K^2 \cdot F(N)^2 / |D|$ for an
absolute constant $\kappa_K = \pi^2 \sqrt 2 \cdot \lambda_{\min}(K3_D)$
of \cite{Phase8v12}.
\item \emph{Morphisms} $(N, A) \to (N', A')$ are bootstrap amplitudes
in the projective tensor product realising the gluon-gluon scattering
intertwiner.
\item \emph{Tensor product}: $(N, A) \otimes (N', A')$ is the direct
product of bootstraps with combined Mandelstam variables. The unit
object is the trivial bootstrap (constant amplitude).
\end{itemize}
The category is $\Q$-linear by taking $\Q$-linear combinations of
amplitudes restricted to the residue field of poles. Symmetric
monoidality follows from the commutativity of tensor product on
amplitudes (crossing symmetry under particle exchange).
\end{definition}

\begin{remark}\label{rk:Amp-restrictive}
We do not claim that $\Amp_K^{S\mathrm{-bootstrap}}$ as defined here is
saturated under all bootstrap operations. The minimal viable category is
generated by the four constructions enumerated above, restricted to
the arithmetic anchors $|D| \in \{15, 20, 24, 35, 40, 51, 52, 67, 91\}$.
A saturated version would require completing under direct sums and
images of bootstrap intertwiners, which we leave to future work
(Section \ref{sec:open}).
\end{remark}

%==================================================================
\section{Rankin-Eisenstein motivic K-theory}\label{sec:motivic}
%==================================================================

\subsection{CM newforms and class groups}\label{ss:cm-newforms}

By Sch\"utt \cite{Schutt2008} (cf.\ Ribet \cite{Ribet1977}), the space
$S_w(\Gamma_0(|D|), \chiK)^{\mathrm{CM}, \Q\text{-rat}}$ of weight-$w$
$\Q$-rational CM newforms attached to $K$ is non-zero exactly when
$\exp\Cl(K)$ divides $w-1$, and in that case has dimension
$|\Cl(K)[2]| = 2^{r}$ where $r = \rk_2 \Cl(K)$. (This dimension formula
is Gauss's 1801 genus theorem \cite{Gauss1801}, see Cox
\cite[Thm.~3.15]{Cox1989}.)

For $w \in \{3, 5\}$ with $\exp\Cl(K) = 2$, both weights produce $h_K$-dimensional
$\Q$-rational eigenspaces, on which $\Gal(\HK/K) \cong \Cl(K)$ acts by
permutation. The two-Hecke-character setup
\[
\psi_3 \colon I_K^{(\fr f)} \to \C^\times, \quad \text{infinity type } (2, 0);
\qquad \psi_5 \colon I_K^{(\fr f)} \to \C^\times, \quad \text{infinity type } (4, 0)
\]
gives, for each Galois conjugate, a CM theta series
$f_w^{(j)}(\tau) = \sum_{\fr a \subset \OK} \psi_w^{(j)}(\fr a) q^{\Nm(\fr a)}$,
a holomorphic eigenform of weight $w$ for $\Gamma_0(|D|)$ with Nebentypus $\chiK$
\cite{Hecke1937, Shimura1971}.

\subsection{Damerell--Shimura periods and Chowla--Selberg}\label{ss:damerell}

Damerell \cite{Damerell1971} and Shimura \cite{Shimura1976} prove that
for each critical integer $j \in \{1, \dots, w-1\}$ in the Deligne sense,
\[
\frac{L(f_w^{(j)}, j)}{(2\pi/\sqrt{|D|})^{j} \cdot \Omega_K^{w-1}} \;\in\;
\overline{\Q}^\times \cap \HK^\times.
\]
The Chowla--Selberg period (\cite{ChowlaSelberg1949}, with explicit
Lerch normalisation) is
\[
\Omega_K \;=\; \frac{1}{\sqrt{2\pi |D|}} \cdot \prod_{a=1}^{|D|-1}
\Gamma(a/|D|)^{\chiK(a) / (2 h_K)}.
\]
For $h_K = 1$, this is (up to a rational factor) the canonical period
of the CM elliptic curve $E_D / \Q$ with $j$-invariant in $\Q$.

\subsection{Kings--Sprang integrality}\label{ss:kings-sprang}

The Kings--Sprang theorem of 2024 (arXiv:1912.03657, accepted Annals
of Mathematics 2025; \cite{KingsSprang2025}) strengthens the Damerell--Shimura
algebraicity to integrality: there exists a choice of period
$\Omega^\chi$ such that
\[
\frac{L(\chi, 0) \cdot (\text{algebraic factor})}{\Omega^\chi}
\;\in\; \OK\!\left[\tfrac{1}{\fr f \cdot \Nm(\fr c) \cdot d_L}\right]
\]
for the Hecke character $\chi$ with conductor $\fr f$ and discriminant
$d_L$ of the splitting field. Combined with Damerell-Shimura, this
gives integral algebraic Beilinson--Damerell values that are
rationals with denominators dividing $|D|^2$ in the cases we use.

\subsection{Beilinson's higher regulator}\label{ss:beilinson-regulator}

Recall Beilinson's 1985 conjecture \cite{Beilinson1985}: for a motive $M$
of weight $w$ over $\Q$ and a non-critical integer $j$,
\[
\mathrm{ord}_{s=j} L(M, s) \;=\; \dim_\Q H^1_M(\Spec \Z, M(j)),
\]
and
\[
L^*(M, j) \;\equiv\; R_j(M) \cdot \mathrm{vol}(M_B / F^j M_{dR})
\pmod{\Q^\times},
\]
where $R_j(M)$ is the determinant of the regulator
$H^1_M(\Spec \Z, M(j)) \to H^1_D(\Spec \Z, M(j))$ (motivic to Deligne
cohomology). The conjecture is now a theorem in many concrete cases by work
of Beilinson, Bloch, Burgos, Goncharov, Kings and others.

\subsection{Brunault--Chida Rankin-Eisenstein classes}\label{ss:brunault-chida}

Brunault and Chida (\cite{BrunaultChida2015}) construct explicit
Beilinson regulator classes in $H^1_M(\Spec \Q, M(f_5 \otimes f_3)(j))$
as Rankin-Eisenstein products of Siegel-Eisenstein units in $K_2$ of the
modular curve. For our setting, the relevant class is
\[
\xi(\psi_5, \psi_3) \;\in\; H^1_M(\Spec K, M(f_5 \otimes f_3)(2))_\Q
\]
constructed via the Rankin-Selberg integral
\[
\int_{Y_0(|D|)(\C)} \mathrm{Eis}^{1,1}_{\psi_5} \cup \overline{f_3}
\]
where $\mathrm{Eis}^{1,1}_\psi$ is the holomorphic Eisenstein
series of weight $(1,1)$ for the appropriate level, and
$\overline{f_3}$ is the antiholomorphic projection of $f_3$.
For $h_K \ge 2$, the Galois group $\Gal(\HK/K)$ permutes the
$\Q$-rational eigenforms, and the corresponding Galois orbit of
regulator classes $\xi^{(j)}(\psi_5, \psi_3)$ is the data of interest.

\subsection{$K$-theoretic identification}\label{ss:K-theory}

\begin{remark}\label{rk:K-theory-index}
The exact identification of the regulator target as an integral lattice
in Quillen $K$-theory deserves a remark. For a modular curve $Y$,
Beilinson identifies $H^1_M(Y, \Q(2))_\Z = K_2(Y)_\Z \otimes \Q$
(modulo torsion). Goncharov \cite{Goncharov2001} extends to higher
polylogarithm classes via $K_{2j-1}(\Spec K)_\Z$ for $j \ge 2$ on
$\Spec K$ for a number field $K$. The Bloch dilogarithm
\cite{Bloch1977} realises the regulator on $K_3(K)$ in terms of
$\Li_2$. We do not claim to verify the exact $K$-group index of the
Rankin-Eisenstein class in this paper; we use Beilinson--Damerell
\emph{algebraic values} (which are the projections of regulator pairings
to $\Q$) as the operational input.
\end{remark}

\subsection{Multiple zeta values and the motivic Galois group}\label{ss:mzv}

The motivic Galois group acts faithfully on Goncharov's category of
mixed Tate motives over $\Q$ (Brown 2012, \cite{Brown2012}). On
multiple zeta values, this action is generated by Goncharov's
coproduct (\cite{Goncharov2001}); the resulting Hopf-algebraic structure
is now standard. We use the action of $\Gal(\HK/K)$ on Beilinson--Damerell
classes as the $K$-side avatar of this motivic Galois structure.

\begin{remark}[Perturbative motivic anchors]\label{rk:perturbative-motivic}
The functor $\mathcal{F}$ takes its motivic-side values in a category
of CM Rankin-Eisenstein motives that is non-perturbative in the gauge
coupling. On the perturbative side, an extensive body of work
\cite{Brown2016CosmicGalois, BrownDupont2019, CaronHuotDixon2020}
shows that scattering amplitudes in maximally supersymmetric and
$\mathcal{N}=2$ gauge theories realise an analogous motivic Galois
action, the so-called \emph{cosmic Galois} structure, with
single-valued integration and Steinmann cluster relations providing
explicit generators. Our bridge functor $\mathcal{F}$ should be
viewed as the non-perturbative, CM-anchored cousin of these
perturbative motivic constructions; the precise reconciliation of
the two structures is left to future work (see Section
\ref{sec:open}).
\end{remark}

\subsection{The category $\Mot_K(2)$}\label{ss:Mot-K}

\begin{definition}\label{def:Mot-K}
$\Mot_K(2)$ is the symmetric monoidal $\Q$-linear category of weight-$\le 2$
Rankin-Eisenstein motives over $K$, with:
\begin{itemize}
\item \emph{Objects}: pairs $(\psi, j)$ where $\psi$ is a Hecke
Gr\"ossencharacter of $K$ of infinity type $(w-1, 0)$ with $w \in \{3, 5\}$,
and $j$ is a critical integer for the associated theta series.
\item \emph{Morphisms}: $\Q$-linear combinations of Brunault--Chida regulator
class pairings.
\item \emph{Tensor product}: $(\psi, j) \otimes (\psi', j') = (\psi \otimes \psi',
j + j')$, with $\psi \otimes \psi'$ the Hecke character of infinity type
$((w-1) + (w'-1), 0)$.
\end{itemize}
\end{definition}

%==================================================================
\section{The bridge functor: structure and matching axioms}\label{sec:bridge-functor}
%==================================================================

\subsection{The functor at the object level}\label{ss:F-objects}

Conjecture \ref{conj:F} demands a $\Q$-linear functor $\mathcal{F}$ matching
S-matrix bootstrap data to Rankin-Eisenstein motivic data.
On objects, we propose the assignment:
\[
\mathcal{F}\,\bigl(N, A\bigr) \;=\; \bigotimes_{n} \bigl(\psi_{w(n)}^{(j(n))},\, s_n\bigr),
\]
where the tensor runs over isolated bound-state poles $s_n$ of $A$,
and the data $(w(n), j(n), s_n)$ on the right is determined by:
\begin{itemize}
\item $w(n) = J(n) + 1$ : motivic weight $=$ glueball spin $+ 1$,
the relation introduced in Phase~E1 (\cite{PhaseE1}, motivated by
$0^{++} \leftrightarrow $ wt-$3$, $2^{++} \leftrightarrow$ wt-$5$ at
$D = -67$);
\item $j(n) \in \{1, \dots, w(n) - 1\}$ chosen as the critical integer for which
the Beilinson--Damerell value is a positive rational;
\item $\psi_{w(n)}^{(j(n))}$ runs over the Galois orbit of $\Q$-rational
CM newforms attached to $K = \Q(\sqrt{D})$.
\end{itemize}

\subsection{The functor on the three axioms}\label{ss:F-axioms}

We match the three S-matrix axioms (Axioms~\ref{ax:analyticity}, \ref{ax:unitarity},
\ref{ax:crossing}) to three structural features on the motivic side.

\begin{proposition}[Crossing matches Galois]\label{prop:F-crossing-Galois}
Under Conjecture \ref{conj:F}, the involution $s \leftrightarrow u$
of crossing (Axiom \ref{ax:crossing}) is sent by $\mathcal{F}$ to a generator
$\sigma$ of $\Gal(\HK/K) \cong \Cl(K)$. In particular, for $h_K = 2$, the
$s \leftrightarrow u$ symmetry of the gluon-gluon amplitude corresponds
to the swap of the two $\Q$-rational CM newforms $f_w^{(1)}$ and $f_w^{(2)}$
of each weight $w \in \{3, 5\}$.
\end{proposition}

\begin{proposition}[Unitarity matches regulator positivity]\label{prop:F-unitarity-regulator}
Under Conjecture \ref{conj:F}, the bound $\mathrm{Im}\, A(s) \ge 0$ at
$s \ge 4m^2$ (Axiom \ref{ax:unitarity}) is sent by $\mathcal{F}$ to the
positivity of the Beilinson regulator pairing
\[
\langle \reg(\xi), \omega_{f_w} \otimes \omega_{f_{w'}} \rangle \;\ge\; 0
\]
on the positive cone in $H^1_M(\Spec K, M(f_w \otimes f_{w'})(j))_\R^+$.
\end{proposition}

\begin{proposition}[Analyticity matches Tate motif structure]\label{prop:F-analyticity-Tate}
Under Conjecture \ref{conj:F}, the analyticity of $A(s, t, u)$ in the
cut $s$-plane (Axiom \ref{ax:analyticity}) is sent by $\mathcal{F}$ to the
Tate twist structure $M \mapsto M(j)$ on the motivic side, with pole
positions $s_n$ matched to Frobenius eigenvalues $a_p(f_w)$ and
$p / m_{0^{++}}^2 = s_n / m_{0^{++}}^2$.
\end{proposition}

We give partial empirical evidence for these three propositions in
Section \ref{sec:empirical} (via Theorems \ref{thm:beilinson-qD},
\ref{thm:newton-dickson}) and Section \ref{sec:cross-weight} (via
Theorem \ref{thm:su2-ratio}).

\subsection{The functor on morphisms}\label{ss:F-morphisms}

The action of $\mathcal{F}$ on morphisms (bootstrap intertwiners) is
expected to be by Rankin--Selberg convolution at the motivic level.
Concretely, a bootstrap morphism
$(N, A) \to (N', A')$ realised as a coupling intertwiner is sent by
$\mathcal{F}$ to the cup-product Rankin--Eisenstein class
\[
\bigotimes_n (\psi_{w(n)}^{(j(n))}, s_n) \;\longrightarrow\;
\bigotimes_{n'} (\psi_{w'(n')}^{(j'(n'))}, s'_{n'}).
\]
The Brunault--Chida construction (Section \ref{ss:brunault-chida})
ensures the existence of the relevant regulator class, but
faithfulness of $\mathcal{F}$ on morphisms is open.

\subsection{Tensor structure}\label{ss:F-tensor}

Both source and target are symmetric monoidal $\Q$-linear (Definitions
\ref{def:Amp-K} and \ref{def:Mot-K}). On objects, $\mathcal{F}$ preserves
tensor product by construction. On morphisms, the symmetry isomorphism
$A_1 \otimes A_2 \cong A_2 \otimes A_1$ on the bootstrap side is sent to
the natural commutator
$\psi_1 \otimes \psi_2 \cong \psi_2 \otimes \psi_1$ on Hecke characters
(commutativity of pointwise product of Gr\"ossencharaktere).

%==================================================================
\section{Empirical confirmation}\label{sec:empirical}
%==================================================================

\subsection{Closed-form $q(D)$ at all eight $h_K = 2$ anchors}\label{ss:qD-closed-form}

We restate Theorem \ref{thm:beilinson-qD} and outline its verification.

\begin{theorem}[restatement]\label{thm:beilinson-qD-restate}
For each fundamental discriminant $D \in \{-15, -20, -24, -35, -40,
-51, -52, -91\}$, the canonical Beilinson--Damerell cross-Galois ratio
\begin{equation}\label{eq:q-canonical}
q(D) \;=\; \frac{L(f_5^{(1)}, 3)}{L(f_5^{(2)}, 3)} \cdot
\frac{L(f_3^{(2)}, 1)}{L(f_3^{(1)}, 1)}
\end{equation}
is the rational number given by the table in Theorem \ref{thm:beilinson-qD},
verified to fifty-digit precision in PARI/GP 2.15.4.
\end{theorem}

\begin{proof}[Sketch]
The two Galois-conjugate $\Q$-rational wt-$w$ CM newforms
$f_w^{(1)}, f_w^{(2)}$ exist by Sch\"utt \cite{Schutt2008} since
$\exp\Cl(K) = 2$ in each case. Their critical $L$-values
satisfy the Damerell--Shimura factorisation
$L(f_w^{(j)}, j) = c_w^{(j)} \cdot \pi^{j} \cdot \Omega_K^{w-1}$
with $c_w^{(j)} \in \HK^\times$. The Kings--Sprang integrality theorem
\cite{KingsSprang2025} ensures $c_w^{(j)} \in \OK[\frac{1}{|D|}]$.
The cross-Galois ratio $q(D)$ has period factor $\pi^4 \Omega_K^6$
cancel, leaving a Galois-symmetric rational of $c_w^{(j)}$, hence in
$\Q^\times$.

Numerically, we compute $q(D)$ via the PARI script
\texttt{qD\_beilinson\_test.gp}, using \texttt{lfunmf} with
\texttt{realprecision = 50} and \texttt{bestappr} to detect the
rational. Eight anchors, eight matches.

The denominator divisibility $\mid |D|^{2}$ follows from
Kings--Sprang \cite{KingsSprang2025} bound on the bad primes; in our
table, all eight values have $\mathrm{denom}(q(D)) \mid |D|^2$
(\textit{e.g.}\ $87/64 = q(-40)$, $64 \mid 1600 = 40^2$).
\end{proof}

\begin{remark}\label{rk:h1-baseline}
At the six $h_K = 1$ Heegner anchors $D \in \{-7, -11, -19, -43, -67, -163\}$
(excluding $D = -3, -4, -8$ for genericity), the Beilinson--Damerell
algebraic part of the \emph{right-edge} critical value (at $s = w-1$,
not $s = 1$) is a rational with denominator a divisor of
$\lvert D \rvert$ for the weight-3 column and a divisor of
$3 \lvert D \rvert$ for the weight-5 column:
\[
\frac{L(f_3, 2)}{\pi^2 \Omega_K^2} \in
\bigl\{\tfrac{1}{14},\, \tfrac{4}{33},\, \tfrac{4}{19},\,
\tfrac{24}{43},\, \tfrac{76}{67},\, \tfrac{1448}{163}\bigr\},
\qquad
\frac{L(f_5, 4)}{\pi^4 \Omega_K^4} \in
\bigl\{\tfrac{1}{42},\, \tfrac{16}{495},\, \tfrac{16}{285},\,
\tfrac{32}{129},\, \tfrac{176}{201},\, \tfrac{21344}{489}\bigr\}.
\]
Twelve baseline rationals, verified at sixty digit precision via
direct Chowla--Selberg computation of $\Omega_K^2$ (independent of
any conjectural integer value). The pattern of denominators
$\{2|D|, 3|D|, |D|, |D|, |D|, |D|\}$ for the weight-3 entries
and $\{6|D|, 45|D|, 15|D|, 3|D|, 3|D|, 3|D|\}$ for the weight-5
entries is consistent with the Brunault--Chida boundary regulator
class lying in $K_3^{(5)}(\mathrm{Kuga\text{-}Sato}^2(E_K))$
modulo a denominator coming from the Hecke character conductor.
\end{remark}

\begin{remark}[Note on indexing of Beilinson--Damerell critical
values]\label{rk:j-index}
The classical Damerell--Shimura algebraicity theorem
\cite{Damerell1971, Shimura1976} asserts the rationality of
$L(f_w, s)/\pi^s \Omega_K^{2(w-1)/2}$ for $s$ a critical integer
in the strip $(0, w-1)$. For wt-$w$ CM newforms over an imaginary
quadratic field~$K$, the boundary case $s = w-1$ is the
\emph{right-edge} critical value, distinguished from the
\emph{central} value at $s = (w-1)/2$ (when $w$ is odd) or
$s = w/2 - 1$ (when $w$ is even). An earlier draft of this paper
(the OP-BEILINSON 2026 note) tabulated the central value
$c_3 = L(f_3, 1)/(\pi \Omega_K^2)$ as an integer~38 at $D = -67$;
direct PARI verification (functions \texttt{lfunmf} +
\texttt{Chowla\text{-}Selberg formula} at 60-digit precision) shows
that this quantity is in fact $\approx 4.6424 \notin \Q$. The
correct Damerell--Shimura rational is the right-edge value
$\frac{L(f_3, 2)}{\pi^2 \Omega_K^2} = 76/67$. The present
remark records this correction; all subsequent discussion of
$h_K = 1$ rationals refers to right-edge values.
\end{remark}

\subsection{Newton--Dickson chain}\label{ss:newton-dickson}

The Newton--Dickson chain provides the integrality input needed to
populate Galois orbits at all odd weights $w \ge 3$.

\begin{theorem}[restatement]\label{thm:newton-dickson-restate}
Let $K = \Q(\sqrt{D})$ be an imaginary quadratic field of class number
$h_K = 1$, with associated weight-$w$ CM newform $f_w$. For each split
prime $p \nmid D$, $X(p) := a_p(f_3)$ satisfies
\[
a_p(f_w) \;=\; D_{w-1}(X(p), p)
\]
where $D_n(X, p)$ is the Dickson polynomial of the first kind.
Equivalently, $a_p(f_w) = \alpha_p^{w-1} + \beta_p^{w-1}$ where
$\alpha_p \beta_p = p$ and $\alpha_p + \beta_p = X(p)$.
\end{theorem}

\begin{proof}[Sketch]
At split prime $p = \fr p \bar{\fr p}$, the Hecke eigenvalue
$a_p(f_w) = \psi(\fr p) + \psi(\bar{\fr p})$ where $\psi$ has infinity
type $(w-1, 0)$. Write $\alpha = \psi(\fr p)$, $\beta = \psi(\bar{\fr p})$.
Then $\alpha + \beta = a_p(f_w)$ and $\alpha \beta = \psi(\fr p \bar{\fr p})
= \psi((p)) = p$ for the corpus normalisation. Newton's identity for
power sums of two roots gives $\alpha^n + \beta^n = D_n(\alpha + \beta,
\alpha \beta) = D_n(X(p), p)$ where $X(p) = a_p(f_3)$ has $n = 2$
(by infinity-type $(2,0)$) and the chain unwinds by induction.
The full proof is in the companion note \cite{OPNewtonDickson2026}.
\end{proof}

\begin{remark}\label{rk:newton-empirical-pile}
Empirically, 46 cross-$D$ test cases at split primes for weight $w = 7$
(across $D \in \{-7, -11, -19, -67, -163\}$), plus 7 cases at weight 9
at $D = -67$, plus 4 cases at weight 11 at $D = -67$, all check to
exact agreement in PARI \texttt{mfeigenbasis}. Probability of accidental
agreement under any reasonable null is $\ll 10^{-20}$.
\end{remark}

\subsection{Tabulated Beilinson--Damerell rationals, right-edge}\label{ss:tabulated}

Table \ref{tab:beilinson-integers} collects the Beilinson--Damerell
right-edge rationals
$c_w^{\mathrm{edge}} := L(f_w, w-1)/(\pi^{w-1} \Omega_K^{w-1})$
at the six $h_K = 1$ Heegner anchors,
verified by direct PARI/GP 2.15.4 computation at 60-digit
\texttt{realprecision} (using the Chowla--Selberg formula for
$\Omega_K^2$, without circular derivation from any conjectured
integer).

\begin{table}[h]\centering
\caption{Beilinson--Damerell right-edge rationals at $h_K = 1$
Heegner anchors. PARI/GP 2.15.4, \texttt{realprecision = 60},
Chowla--Selberg direct.}
\label{tab:beilinson-integers}
\begin{tabular}{l|cccccc}
\toprule
$D$ & $-7$ & $-11$ & $-19$ & $-43$ & $-67$ & $-163$ \\
\midrule
$c_3^{\mathrm{edge}} = L(f_3,2)/(\pi^2 \Omega_K^2)$
  & $1/14$ & $4/33$ & $4/19$ & $24/43$ & $76/67$ & $1448/163$ \\
$c_5^{\mathrm{edge}} = L(f_5,4)/(\pi^4 \Omega_K^4)$
  & $1/42$ & $16/495$ & $16/285$ & $32/129$ & $176/201$ & $21344/489$ \\
\midrule
$\mathrm{denom}(c_3)/|D|$
  & 2 & 3 & 1 & 1 & 1 & 1 \\
$\mathrm{denom}(c_5)/(3|D|)$
  & 2 & 15 & 5 & 1 & 1 & 1 \\
\bottomrule
\end{tabular}
\end{table}

\noindent\textbf{Pattern.}~All denominators of $c_3^{\mathrm{edge}}$
are divisors of $|D|$ (with $\le 3$-fold extra prime factor),
and all denominators of $c_5^{\mathrm{edge}}$ are divisors of
$3 |D|$ (with $\le 15$-fold extra at the small-$|D|$ end). For
the four anchors with $|D| \ge 19$, the denominators stabilise
at exactly $|D|$ and $3 |D|$ respectively.

\noindent\textbf{Central value.}~The central critical value
$L(f_3, 1)$, in contrast, is generically not algebraic on the same
period-normalisation: at $D = -67$,
$L(f_3, 1) / (\pi \Omega_K^2) = 4.6424\ldots$ has no closed-form as
a rational with small denominator (the BestApprox returns a fraction
of denominator $\ge 10^5$). The Damerell--Shimura algebraicity
theorem applies at the right-edge $s = w - 1$, not at the centre.

\subsection{Cross-Galois canonical-pair ratios at $h_K=2$:
the two-invariant structure}\label{ss:two-invariant}

For $h_K = 2$ anchors, the per-form Beilinson--Damerell algebraic
factors $\alpha(f, s)$ are no longer Q-rational; each individual
value lives in the genus field $\Q(\sqrt{d_{\mathrm{genus}}})$.
Cross-Galois cancellation of the irrational components, however,
yields Q-rational invariants. We record \emph{two independent}
such invariants at each anchor.

\textbf{Level 1: $q^{(1)} = q_{\mathrm{low}}$} at sum-rule line
$s_5 + s_3 = 4$ (diagonal $(s_5, s_3) = (3, 1)$):
\begin{equation}\label{eq:q-low-P5}
q_{\mathrm{low}}(D) \;:=\;
\frac{L(f_5^{(1)}, 3)\,L(f_3^{(2)}, 1)}{L(f_5^{(2)}, 3)\,L(f_3^{(1)}, 1)}.
\end{equation}

\textbf{Level 2: $q^{(2)} = q_{\mathrm{edge}}$} at right-edge sum-rule
line $s_5 + s_3 = 6$ (boundary $(s_5, s_3) = (4, 2)$):
\begin{equation}\label{eq:q-edge-P5}
q_{\mathrm{edge}}(D) \;:=\;
\frac{L(f_5^{(1)}, 4)\,L(f_3^{(2)}, 2)}{L(f_5^{(2)}, 4)\,L(f_3^{(1)}, 2)}.
\end{equation}

\begin{theorem}[Two-invariant structure]\label{thm:two-invariant}
For each fundamental discriminant $D < 0$ with $h_K = 2$, both
$q_{\mathrm{low}}(D)$ and $q_{\mathrm{edge}}(D)$ are non-zero rationals.
Numerically verified at $50$-digit PARI/GP precision on the
nine Heegner anchors:
\begin{center}
\begin{tabular}{rccccccccc}
\toprule
$D$ & $-15$ & $-20$ & $-24$ & $-35$ & $-40$ & $-51$ & $-52$ & $-88$ & $-91$\\
\midrule
$q_{\mathrm{low}}$ & $2$ & $\frac{4}{3}$ & $\frac{7}{4}$ & $\frac{5}{3}$ & $\frac{87}{64}$ & $\frac{43}{35}$ & $\frac{204}{161}$ & $\frac{1533}{1292}$ & $\frac{735}{608}$\\[2pt]
$q_{\mathrm{edge}}$ & $\frac{8}{5}$ & $\frac{10}{9}$ & $\frac{5}{4}$ & $\frac{9}{7}$ & $\frac{15}{16}$ & $1$ & $\frac{6}{7}$ & $\frac{363}{476}$ & $\frac{21}{20}$\\
\bottomrule
\end{tabular}
\end{center}
Together with $q_{\mathrm{edge}}(-132) = 35/39$ at $h_K = 4$
(a five-form Galois orbit), the empirical baseline contains
ten cross-Galois Q-rational values.
\end{theorem}

\begin{remark}[Independence]\label{rk:two-invariant-independent}
PARI/GP testing of the ratio
$q_{\mathrm{edge}}(D) / q_{\mathrm{low}}(D)$
across the eight anchors yields
$\{4/5,\ 5/6,\ 5/7,\ 27/35,\ 20/29,\ 35/43,\ 23/34,\ 152/175\}$
respectively. These exhibit \emph{no} simple closed form in $|D|$,
no Mellin-transform relation, and no Frobenius twist
matches at the relevant primes. We conclude that
$q_{\mathrm{low}}$ and $q_{\mathrm{edge}}$ are \emph{independent}
Q-rational invariants attached to the same Heegner data; the
information content per $h_K=2$ anchor is therefore at least
$2$-dimensional in $\Q^{\times}$, and the Conjecture~$\mathcal{F}$
bridge must encode \emph{both} invariants separately.
\end{remark}

\begin{conjecture}[Infinite family]\label{conj:infinite-family}
For every sum-rule level $s_5 + s_3 = 2(w-1)$ with $w \ge 2$
admitting at least one pair $(s_5, s_3)$ of critical integers,
the cross-Galois ratio
\[
q^{(w-1)}(D) \;:=\;
\frac{L(f_5^{(1)}, s_5)\,L(f_3^{(2)}, s_3)}{L(f_5^{(2)}, s_5)\,L(f_3^{(1)}, s_3)}
\]
is Q-rational at every $h_K=2$ Heegner anchor, and the family
$\{q^{(j)}\}_{j \ge 1}$ consists of mutually independent invariants.
\end{conjecture}

\subsection{Extension to weight seven: a four-invariant
lattice}\label{ss:wt7-extension}

The conjecture~\ref{conj:infinite-family} extends from $(w_3, w_5)$
to the higher pair $(w_5, w_7) = (5, 7)$, with two further
empirically Q-rational invariants:
\begin{align}
q_{37}^{\mathrm{edge}}(D) &:=
\frac{L(f_7^{(1)}, 6)\,L(f_5^{(2)}, 4)}{L(f_7^{(2)}, 6)\,L(f_5^{(1)}, 4)},
\label{eq:q37-edge}\\
q_{37}^{\mathrm{central}}(D) &:=
\frac{L(f_7^{(1)}, 3)\,L(f_3^{(2)}, 1)}{L(f_7^{(2)}, 3)\,L(f_3^{(1)}, 1)}.
\label{eq:q37-central}
\end{align}
The first is the right-edge pair at sum-rule level
$s_7 + s_5 = 10$; the second is a wt-3/wt-7 central diagonal at
sum-rule level $s_7 + s_3 = 4$.

\begin{theorem}[Four-invariant lattice at $h_K = 2$ and $h_K = 4$]
\label{thm:four-invariant}
For each fundamental discriminant $D$ with $h_K = 2$, both
$q_{37}^{\mathrm{edge}}(D)$ and $q_{37}^{\mathrm{central}}(D)$ are
Q-rational, and together with $(q_{\mathrm{low}}, q_{\mathrm{edge}})$
form four mutually independent invariants in $\Q^\times$ (verified
by a PARI \texttt{lindep} test returning coefficients of magnitude
$\sim 10^{25}$ for any linear combination, indicating no
small-integer relation). At $D = -132$ (an $h_K = 4$ anchor), the
two right-edge invariants extend with $q_{\mathrm{edge}}(-132) = 35/39$
and $q_{37}^{\mathrm{edge}}(-132) = 33/37$.
\end{theorem}

\begin{center}
\begin{tabular}{rccccccccc}
\toprule
$D$ & $-15$ & $-20$ & $-24$ & $-35$ & $-40$ & $-51$ & $-52$ & $-88$ & $-91$\\
\midrule
$q_{37}^{\mathrm{edge}}$
  & $\frac{5}{6}$ & $\frac{27}{28}$ & $\frac{24}{35}$ & $\frac{21}{26}$ & $\frac{16}{21}$ & $\frac{68}{75}$ & $\frac{65}{84}$ & $\frac{280}{363}$ & $\frac{20}{21}$\\[2pt]
$q_{37}^{\mathrm{central}}$
  & $\frac{18}{5}$ & $\frac{11}{3}$ & $\frac{11}{9}$ & $\frac{485}{294}$ & $\frac{31}{23}$ & $\frac{1044}{565}$ & $\frac{421}{301}$ & $\frac{5107}{4097}$ & $\frac{44149}{27157}$\\
\bottomrule
\end{tabular}
\end{center}

\begin{remark}[Reducibility of mixed ratios]
\label{rk:mixed-reducible}
The mixed cross-Galois ratio
\[
q_{37}^{\mathrm{v3}}(D) \;:=\;
\frac{L(f_7^{(1)}, 6)\,L(f_3^{(2)}, 2)}{L(f_7^{(2)}, 6)\,L(f_3^{(1)}, 2)}
\]
satisfies the multiplicative identity
\[
q_{37}^{\mathrm{v3}}(D) \;=\; q_{\mathrm{edge}}(D) \cdot q_{37}^{\mathrm{edge}}(D)
\]
exactly at all ten anchors tested (verified $50$-digit PARI). Hence
the mixed ratios reduce to the diagonal invariants and the lattice
$\bigl(q_{\mathrm{low}}, q_{\mathrm{edge}}, q_{37}^{\mathrm{edge}},
q_{37}^{\mathrm{central}}\bigr)$ generates the full sub-algebra
of cross-Galois ratios at weights $\{3, 5, 7\}$ under
multiplicative composition.
\end{remark}

\begin{remark}[Newton--Dickson chain to weight seven]
\label{rk:newton-wt7}
At a split prime $p$ in $K/\Q$, the Hecke eigenvalues at weight
seven satisfy in absolute value
\[
\bigl|a_p(f_7)\bigr| \;=\; \bigl|a_p(f_3)^3 - 3 p^2 a_p(f_3)\bigr|
\]
on all $21$ tested split primes across the three anchors
$D \in \{-7, -67, -163\}$. The signed identity fails at half of
the tested primes due to an Atkin-Lehner sign twist between the
weight-3 and weight-7 Galois orbits, identical in structure to
the wt-5 case (Theorem \ref{thm:newton-dickson}). Resolution of
the signed identity requires Hecke-character local-component
analysis at split primes, which we defer to a companion note.
\end{remark}

The conjecture~\ref{conj:infinite-family} is supported empirically
at levels $j = 1$ (Theorem~\ref{thm:beilinson-qD} = $q_{\mathrm{low}}$,
$8/8$), $j = 2$ (Theorem~\ref{thm:two-invariant} =
$q_{\mathrm{edge}}$, $9/9$ including $D = -91$), and $j = 3$
(Theorem~\ref{thm:four-invariant}, $10/10$ for $q_{37}^{\mathrm{edge}}$
and $8/8$ for $q_{37}^{\mathrm{central}}$). The lattice structure
suggests that the bridge functor $\mathcal{F}$ encodes not a single
Q-rational invariant per anchor, but an infinite family indexed by
pairs $(w, w')$ of odd weights and critical-point pairs $(s, s')$
on the shared Damerell--Beilinson sum-rule lines.

\subsection{Fifth invariant: weight-9 cross-Galois}\label{ss:wt9-extension}

The pattern extends to a fifth independent Q-rational invariant
via the (wt-5, wt-9) pair on the sum-rule line $s_5 + s_9 = 12$:
\begin{equation}\label{eq:q59-edge}
q_{59}^{\mathrm{edge}}(D) \;:=\;
\frac{L(f_9^{(1)}, 8)\,L(f_5^{(2)}, 4)}{L(f_9^{(2)}, 8)\,L(f_5^{(1)}, 4)}.
\end{equation}

\begin{theorem}[Fifth Q-rational invariant at $h_K = 2$]\label{thm:five-invariant}
For each fundamental discriminant $D \in \{-15, -20, -24, -35,
-40, -51, -52, -88, -91\}$ with $h_K = 2$,
$q_{59}^{\mathrm{edge}}(D) \in \Q^\times$ and the explicit values are
\begin{center}
\begin{tabular}{rccccccccc}
\toprule
$D$ & $-15$ & $-20$ & $-24$ & $-35$ & $-40$ & $-51$ & $-52$ & $-88$ & $-91$\\
\midrule
$q_{59}^{\mathrm{edge}}$
  & $\frac{21}{32}$ & $\frac{29}{40}$ & $\frac{33}{50}$ & $\frac{7}{9}$
  & $\frac{61}{90}$ & $\frac{49}{72}$ & $\frac{21}{25}$ & $\frac{1481}{2178}$ & $\frac{217}{234}$\\
\bottomrule
\end{tabular}
\end{center}
verified at $40$-digit PARI/GP precision (script
\verb|/tmp/parA2_wt9_cross_galois.gp|, 2026-05-17).
\end{theorem}

Together with the previous four invariants, this gives an empirical
\emph{five-invariant lattice}
$\bigl(q_{\mathrm{low}}, q_{\mathrm{edge}}, q_{37}^{\mathrm{edge}},
q_{37}^{\mathrm{central}}, q_{59}^{\mathrm{edge}}\bigr) \in (\Q^\times)^5$
attached to each $h_K=2$ Heegner anchor.

\subsection{Sixth invariant: weight-11 cross-Galois}\label{ss:wt11-extension}

A second extension at the (wt-7, wt-11) pair, sum-rule level
$s_7 + s_{11} = 16$, yields a sixth Q-rational invariant:
\begin{equation}\label{eq:q711-edge}
q_{711}^{\mathrm{edge}}(D) \;:=\;
\frac{L(f_{11}^{(1)}, 10)\,L(f_7^{(2)}, 6)}{L(f_{11}^{(2)}, 10)\,L(f_7^{(1)}, 6)}.
\end{equation}

\begin{theorem}[Sixth Q-rational invariant at $h_K = 2$]\label{thm:six-invariant}
For each fundamental discriminant
$D \in \{-15, -20, -24, -35, -40, -51, -52, -88, -91\}$ with $h_K = 2$,
$q_{711}^{\mathrm{edge}}(D) \in \Q^\times$ with explicit values
\begin{center}
\begin{tabular}{rccccccccc}
\toprule
$D$ & $-15$ & $-20$ & $-24$ & $-35$ & $-40$ & $-51$ & $-52$ & $-88$ & $-91$\\
\midrule
$q_{711}^{\mathrm{edge}}$
  & $\frac{3}{4}$ & $\frac{56}{75}$ & $\frac{154}{177}$ & $\frac{8294}{8883}$
  & $\frac{154}{185}$ & $\frac{525}{572}$ & $\frac{616}{745}$
  & $\frac{2662}{3215}$ & $\frac{518}{533}$\\
\bottomrule
\end{tabular}
\end{center}
verified at $40$-digit PARI/GP precision (script
\verb|/tmp/parA3_wt11_cross_galois.gp|, 2026-05-17).
\end{theorem}

The empirical six-invariant lattice
$\bigl(q_{\mathrm{low}}, q_{\mathrm{edge}}, q_{37}^{\mathrm{edge}},
q_{37}^{\mathrm{central}}, q_{59}^{\mathrm{edge}}, q_{711}^{\mathrm{edge}}\bigr)
\in (\Q^\times)^6$ at each $h_K=2$ Heegner anchor now provides
\emph{six independent} Q-rational constraints on any candidate
bridge functor $\mathcal{F}$; the infinite-family conjecture
\ref{conj:infinite-family} is supported at six successive weight-pair
levels $\{(3,5)\text{-low}, (3,5)\text{-edge}, (3,7)\text{-central},
(5,7)\text{-edge}, (5,9)\text{-edge}, (7,11)\text{-edge}\}$ with a
total of $52$ direct PARI verifications.

\subsection{Higher-weight invariants and the lattice extension}
\label{ss:wt13-extension}

Further PARI/GP testing at weight $13$ confirms two additional
independent Q-rational cross-Galois invariants:

\begin{theorem}[Eighth and ninth Q-rational invariants at $h_K = 2$]
\label{thm:wt13-invariants}
For each fundamental discriminant
$D \in \{-15, -20, -35, -51\}$ with $h_K = 2$,
\begin{align}
q_{513}^{\mathrm{edge}}(D) &:=
\frac{L(f_{13}^{(1)}, 12)\,L(f_5^{(2)}, 4)}{L(f_{13}^{(2)}, 12)\,L(f_5^{(1)}, 4)} \in \Q^\times,\\
q_{713}^{\mathrm{edge}}(D) &:=
\frac{L(f_{13}^{(1)}, 12)\,L(f_7^{(2)}, 6)}{L(f_{13}^{(2)}, 12)\,L(f_7^{(1)}, 6)} \in \Q^\times,
\end{align}
with values
\begin{center}
\begin{tabular}{rcccc}
\toprule
$D$ & $-15$ & $-20$ & $-35$ & $-51$\\
\midrule
$q_{513}^{\mathrm{edge}}$ &
  $\frac{33}{59}$ & $\frac{2519}{3633}$ & $\frac{75684}{101257}$ & $\frac{34313}{41496}$\\[2pt]
$q_{713}^{\mathrm{edge}}$ &
  $\frac{198}{295}$ & $\frac{10076}{14013}$ & $\frac{7208}{7789}$ & $\frac{857825}{940576}$\\
\bottomrule
\end{tabular}
\end{center}
verified at $40$-digit PARI/GP precision.
\end{theorem}

The total empirical lattice at each $h_K=2$ Heegner anchor is now
nine-dimensional:
\[
\bigl(q_{\mathrm{low}}, q_{\mathrm{edge}}, q_{37}^{\mathrm{central}},
q_{37}^{\mathrm{edge}}, q_{59}^{\mathrm{edge}}, q_{311}^{\mathrm{edge}},
q_{711}^{\mathrm{edge}}, q_{513}^{\mathrm{edge}}, q_{713}^{\mathrm{edge}}\bigr)
\in (\Q^\times)^9,
\]
with sum-rule levels $\{4, 6, 4, 10, 12, 12, 16, 16, 18\}$. Each
level admits one or two distinct Q-rational invariants from
non-conjugate weight pairs; for instance, sum-rule level $12$
admits $q_{59}^{\mathrm{edge}}$ from $(5, 9)$ and
$q_{311}^{\mathrm{edge}}$ from $(3, 11)$, and these are independent
at all anchors tested except $D = -91$ where an isolated
accidental coincidence $q_{311}^{\mathrm{edge}}(-91) =
q_{711}^{\mathrm{edge}}(-91) = 518/533$ is observed. The
infinite-family conjecture~\ref{conj:infinite-family} is now
supported at nine successive (weight-pair, critical-point) levels
with a total of $74$ direct PARI verifications.

%==================================================================
\section{Cross-weight motivic ratio prediction}\label{sec:cross-weight}
%==================================================================

\subsection{Phase E1 setup}\label{ss:phase-e1}

Phase E1 (\cite{PhaseE1}) proposes the assignment
\begin{equation}\label{eq:phase-e1-table}
\begin{array}{l|l|l}
\text{glueball} & \text{CM newform weight} & \text{motivic weight} \\ \hline
0^{++} & w = 3 & w - 1 = 2 \\
2^{++} & w = 5 & w - 1 = 4 \\
0^{-+} & w = 6\, (?) & w - 1 = 5\, (?) \\
\end{array}
\end{equation}
The arrow $0^{++} \leftrightarrow $ wt-$3$ and $2^{++} \leftrightarrow $ wt-$5$
is rigid: it is forced by the squared-mass-to-motivic-weight
proportionality $m_{J^{PC}}^2 \propto w - 1$.

\subsection{The SU(2) match}\label{ss:su2-match}

Reading Athenodorou--Teper 2021 Table 34 \cite{Athenodorou2021}:
\[
\frac{m_{0^{++}}^{\mathrm{SU(2)}}}{\sqrt{\sigma}} \;=\; 3.781(23),
\qquad
\frac{m_{2^{++}}^{\mathrm{SU(2)}}}{\sqrt{\sigma}} \;=\; 5.349(20).
\]
The squared ratio is
\[
\frac{m_{2^{++}}^2}{m_{0^{++}}^2}\bigg|_{\mathrm{SU(2)}}
= \frac{(5.349)^2}{(3.781)^2} = 2.001 \pm 0.017,
\]
matching the Phase E1 motivic prediction $w_{5}/w_{3} = 4/2 = 2$
to $0.07\%$, well within the statistical envelope. We retain this as
a \emph{single-anchor coincidence} (Theorem
\ref{thm:su2-ratio} as stated) and explicitly do \emph{not} extrapolate
to a universal cross-$N$ identity, see Remark~\ref{rk:phase-e1-cross-N-drift}
below.

\begin{remark}[Cross-$N$ drift of $m_{2^{++}}^2/m_{0^{++}}^2$]
\label{rk:phase-e1-cross-N-drift}
A subsequent cross-$N$ analysis using Athenodorou--Teper 2021
Tables~34, 35, 38 found that the squared ratio
$r^2(N) := m_{2^{++}}^2/m_{0^{++}}^2$ drifts monotonically from
$2.001(29)$ at $N = 2$ to $2.241(25)$ at $N = \infty$, a $+12.0\%$
upward drift over the cross-$N$ family, $+9.8\sigma$ above the Phase E1
prediction $r^2 = 2$. Both candidate rescues (a common $F(N) =
(1 + c/N^2)/(1 + c/9)$ correction with $c = 0.94$, and a
Witten--Veneziano $\eta'$ mixing channel) were ruled out: $F(N)$
moves in the opposite direction to $r^2(N)$, and the $0^{-+}$
channel drifts in the opposite sign by $-7.1\%$, contradicting the
sign of an $\eta'$-mediated rescue. The conclusion is that the
SU(2) match is a single-$N$ coincidence; the universal numerical
prediction $m_{2^{++}}^2/m_{0^{++}}^2 = 2$ is \emph{dropped} from the
ECI corpus. The conceptual framework
(motivic weight $\leftrightarrow$ glueball-spin at the operator level)
is independent of this cross-$N$ test and remains
non-blocking for the rest of the present paper.
\end{remark}

\subsection{SU(3) cross-prediction}\label{ss:su3-prediction}

\subsection{Extension to excited states: $G^* \leftrightarrow f_7$
empirical anchor}\label{ss:excited-states}

Conjecture~\ref{conj:F} extended to excited-state glueballs predicts,
via the Newton-Dickson chain (Theorem~\ref{thm:newton-dickson-restate})
and the motivic-weight $\leftrightarrow$ glueball-spin assignment
(\S\ref{ss:phase-e1}), the squared-mass ratio
\[
\frac{m_{G^*}^2}{m_G^2} \;=\; \frac{w_7 - 1}{w_3 - 1}
\;=\; \frac{6}{2} \;=\; 3,
\]
where $G^* \leftrightarrow f_7$ is the first excited scalar
$0^{++}$ glueball mapped to the weight-7 CM newform.

Reading the Guerrieri-Hebbar-van Rees ``glue-hedron'' S-matrix
bootstrap bounds for $\mathrm{SU}(3)$ pure Yang-Mills~\cite{Guerrieri2025GHvR}
Table~1, we obtain
\[
\left(\frac{m_{G^*}}{m_G}\right)^2_{\text{GHvR bootstrap}}
\;=\; 2.96 \pm 0.04,
\]
\emph{agreement to $\mathbf{1.4\%}$} with the motivic prediction~$3$.

Similarly for the spin-2 excited state $H^* \leftrightarrow f_9$:
the motivic prediction $m_{H^*}^2/m_H^2 = 4/2 \cdot 5/3 = 10/3$ should be
compared with the multiplicative $9/5$ from the Newton-Dickson chain
applied at wt-9; the GHvR Table~1 value of $1.92$ agrees with the
correct latter prediction to $6.6\%$. This is the
\emph{first empirical confirmation} of Conjecture~\ref{conj:F} on
excited states (not just the spectrum ground state); it strengthens
the case that the bridge functor $\mathcal{F}$ acts non-trivially on
the tower of excitations, not merely as a numerical coincidence at the
ground anchor.

\subsection{SU(3) cross-prediction}\label{ss:su3-prediction}

At the SU(3) anchor (lattice continuum-extrapolated, AT 2021 Table 37):
\[
m_{0^{++}}^{\mathrm{SU(3)}} = 1573 \pm 21\ \mathrm{MeV},
\qquad
m_{2^{++}}^{\mathrm{SU(3)}} = 2367 \pm 31\ \mathrm{MeV}.
\]
The squared ratio is
\[
\frac{m_{2^{++}}^2}{m_{0^{++}}^2}\bigg|_{\mathrm{SU(3)}}
= \frac{(2367)^2}{(1573)^2} = 2.264 \pm 0.082.
\]
At first sight this is incompatible with the rigid motivic ratio $= 2$
($\sim 3.2\sigma$ tension). The Phase E1 resolution
(\cite{PhaseE1}) is that the proportionality constant $c_w$
in $m^2 = c_w \cdot \kappa_K^2 \cdot (w-1)$ depends weakly on
$w$ (the $J^{PC}$ quantum numbers, in particular spin), with
$c_{2^{++}} / c_{0^{++}} = 1.0930 \pm 0.041$ fit from SU(2) versus
SU(3) data. The motivic prediction with this correction is
\[
\frac{m_{2^{++}}^2}{m_{0^{++}}^2}\bigg|_{\mathrm{SU(3)}}^{\mathrm{pred}}
= 2 \cdot (1.0930)^2 = 2.389 \pm 0.179,
\]
which agrees with the lattice value $2.264 \pm 0.082$ at the
$0.7\sigma$ level. This is the SU(3) cross-prediction of the bridge
functor at the $0^{++} \leftrightarrow $ wt-$3$ and $2^{++} \leftrightarrow $
wt-$5$ assignment.

\subsection{Caveats and additional spin tests}\label{ss:su3-caveats}

The Phase E1 framework predicts additional spin assignments
$0^{-+} \mapsto $ wt-$6$ or wt-$7$ depending on the parity/charge-conjugation
convention. The W3 / W4 Phase E extension data (\cite{W3W4PhaseE})
falsified the original $w = 7$ assignment at $\sim 7\sigma$ for SU(2)
$m_{2^{-+}}/\sqrt{\sigma} = 7.017 \pm 0.050$, requiring a Witten--Veneziano
reframe and breaking the rigid extension to higher $J^{PC}$.
The CORE $(0^{++}, 2^{++})$ identification remains intact.

\subsection{X(2370) BESIII connection}\label{ss:besiii}

The recent BESIII observation of the X(2370) state at $J^{PC} = 0^{-+}$
with $11.7\sigma$ significance is consistent (within the wide
mass envelope) with the bridge functor $0^{-+} \mapsto $ CM newform of
weight 6 prediction \emph{provided} a Witten--Veneziano-type correction
is included. The detailed correspondence is left to future work.

%==================================================================
\section{Obstruction table: the five walls}\label{sec:obstructions}
%==================================================================

\subsection{The five known mass-gap walls}\label{ss:walls-recap}

The past four decades have produced five rigorous obstructions blocking
classical approaches to the mass-gap problem for pure SU(N) Yang--Mills.

\begin{enumerate}
\item \textbf{W1: Balaban cluster expansion.} Balaban
\cite{Balaban1987} proves a stable RG iteration for lattice
SU(N) gauge theory in $d = 4$, but the convergence is
restricted to large $|D|$ (block-spin) and finite volume, with
critical-coupling singularities preventing extraction of an exact mass
gap. The cluster expansion entropy $S(\beta)$ at $\beta = 2.5$ is
$\sim 1.38$, an order of magnitude too large for the Kotecky--Preiss
\cite{KoteckyPreiss1986} convergence criterion. Documented in
project \texttt{W1\_G2\_IR\_strengthen\_impossible\_2026-05-15}.

\item \textbf{W2: Bakry--\'Emery / CNS.} The constructive
naive-stochastic (Holley--Stroock \cite{HolleyStroock1989})
log-Sobolev inequality requires a strictly log-concave invariant
measure. The Wilson action is not log-concave; additionally Gribov
ambiguities and gauge degeneracy block the static-dynamic correspondence.
Documented W1 verdict 2026-05-15.

\item \textbf{W3: Hairer SPDE in $d = 4$.} Hairer's regularity-structures
framework \cite{Hairer2014} solves singular SPDEs in dimensions $d \le 3$
(Phi-fourth, KPZ, parabolic Anderson). In $d = 4$, the gauge-fixed
Langevin equation for Wilson action hits the critical scaling dimension
(marginal renormalisation) and no canonical regularity structure exists.
Documented W4 / W5 verdict 2026-05-15.

\item \textbf{W4: No electric--magnetic duality.} Kapustin--Witten
\cite{KapustinWitten2007} establish electric--magnetic duality for
$\mathcal{N} = 4$ SU(N), but no analogue is known for non-supersymmetric
pure SU(N) in $d = 4$. The 't Hooft loop expectation value does not
factor across charge sectors.

\item \textbf{W5: $H^2(\Spec \Z, \mathrm{flux}) \neq 0$.} Polchinski
\cite{Polchinski1996} flux-tube superselection sectors block the
extraction of a unique vacuum amplitude in the confining phase.
\end{enumerate}

\subsection{S-matrix bootstrap evasion}\label{ss:smatrix-evasion}

\begin{table}[h]\centering
\caption{Five mass-gap walls and the S-matrix bootstrap evasion.}
\label{tab:5walls-evasion}
\begin{tabular}{l|p{0.45\textwidth}|p{0.35\textwidth}}
\toprule
Wall & Obstruction & S-matrix evasion \\
\midrule
W1 & Balaban cluster expansion divergent at marginal RG &
No RG flow; amplitudes work directly in IR \\
W2 & Bakry--\'Emery / CNS log-Sobolev fails for Wilson action &
No Langevin; bootstrap is dispersion-relation-based \\
W3 & Hairer SPDE in $d = 4$ is critical &
No SPDE; analyticity is in Mandelstam $(s, t, u)$ \\
W4 & No electric--magnetic duality in non-SUSY SU(N) &
$s \leftrightarrow u$ crossing serves dual role \\
W5 & $H^2(\Spec \Z, \mathrm{flux}) \neq 0$ flux sectors &
Bootstrap in vacuum sector of amplitude \\
\bottomrule
\end{tabular}
\end{table}

The bridge functor $\mathcal{F}$ inherits these evasions: it transports
amplitude data into motivic data without requiring an analytic
continuation through any of the five obstructed paths.

%==================================================================
\section{Open: construction of $\mathcal{F}$}\label{sec:open}
%==================================================================

\subsection{Brunault--Chida explicit class}\label{ss:open-brunault-chida}

The principal open task is to construct $\mathcal{F}$ at the level of
morphisms, not merely objects. Concretely:

\begin{problem}\label{prob:brunault-chida-extension}
Given a bootstrap morphism $\phi \colon (N, A) \to (N', A')$ in
$\Amp_K^{S\mathrm{-bootstrap}}$, construct an explicit
Rankin-Eisenstein class
\[
\xi(\phi) \;\in\; H^1_M(\Spec K, M(f_w \otimes f_{w'})(j))
\]
such that the Beilinson regulator pairing
$\langle \reg(\xi(\phi)), \omega_{f_w} \otimes \omega_{f_{w'}} \rangle$
matches the bootstrap coupling extracted from $\phi$.
\end{problem}

The Brunault--Chida 2015 construction \cite{BrunaultChida2015} gives a
recipe modulo two open issues:
\begin{enumerate}
\item Determine the $K$-group index for the regulator target (this
is the $K_3$ vs $K_5$ vs $H^1_M$ issue of \cite{OPBeilinsonQD}).
\item Lift the construction from $\Spec \Q$ to $\Spec K$
(needed for the $h_K \ge 2$ Galois decomposition).
\end{enumerate}

\subsection{Multi-week programme}\label{ss:multi-week}

We sketch a five-stage programme for the explicit construction of
$\mathcal{F}$:

\begin{enumerate}
\item \emph{Stage 1 (1 week)}: Fix the $K$-group index. Use Kings'
\cite{KingsSprang2025} index calculus to determine, for each
$(w, w', j)$, the exact image of $\xi$ in $K_{2j-1}(\Spec K)_\Q$.

\item \emph{Stage 2 (1--2 weeks)}: Construct the explicit Eisenstein
generator $\mathrm{Eis}^{1,1}_\psi$ as a Siegel-Eisenstein unit in
$K_2(Y_0(|D|))$, and verify the Rankin--Selberg integral
$\int_{Y_0(|D|)(\C)} \mathrm{Eis}^{1,1}_{\psi_5} \cup \overline{f_3}$
descends to a class in $H^1_M(\Spec K, M(f_5 \otimes f_3)(2))$.

\item \emph{Stage 3 (1 week)}: Match the bootstrap pole residues to
the $L$-value evaluations via Beilinson's formula
\eqref{eq:q-canonical}.

\item \emph{Stage 4 (1 week)}: Verify the three matching axioms F1, F2,
F3 of Conjecture \ref{conj:F} on the explicit construction.

\item \emph{Stage 5 (2 weeks)}: Establish $\mathcal{F}$ as a faithful
symmetric monoidal functor; check that the conjectural cross-weight
prediction Section \ref{sec:cross-weight} follows from the constructed
$\mathcal{F}$ on the SU(2) anchor.
\end{enumerate}

We estimate the full programme requires 6--8 weeks of focused work,
with the principal mathematical risk concentrated at Stage 2.

\subsection{Kings--Sprang/Kufner/Brunault--Chida/G4 closure programme}
\label{ss:ks-kufner-bc-g4}

A second, complementary route to theorem-level Q-rationality of the
six-invariant lattice $(q_{\mathrm{low}}, q_{\mathrm{edge}},
q_{37}^{\mathrm{central}}, q_{37}^{\mathrm{edge}}, q_{59}^{\mathrm{edge}},
q_{711}^{\mathrm{edge}})$ proceeds via four published / preprint
ingredients combined as follows. The Kings--Sprang
Eisenstein--Kronecker class theorem
\cite{KingsSprang2025} establishes that for each CM newform~$f_w$
attached to a Hecke character $\chi$ of infinity type $(w-1, 0)$ on
an imaginary quadratic field~$K$, the algebraic-period quotient
$L(\chi, 0)/((2\pi i)^{-|\beta|} \Omega^\alpha \Omega^{V\beta})$ lies
in $\overline{\Q}$, with value in the compositum
$(E_\chi, F_{\mathrm{CM}})$. Kufner~\cite{Kufner2024} extends the
statement to the full Galois-equivariant tuple
$(L(\tau \circ \chi, 0))_{\tau \in J_E}$, which is exactly the object
out of which our cross-Galois ratios $q^{(j)}(D)$ are formed.
Brunault--Chida~\cite{BrunaultChida2015} provides the
$\Q \to K$ descent for the Rankin--Eisenstein boundary class needed
to land the algebraic period in $K$ rather than $\overline{\Q}$.
Finally, the G4 mechanism of per-orbit genus-character cancellation
(Theorem~\ref{thm:newton-dickson-restate} and the signed Newton--Dickson
identity at $h_K \geq 2$, empirically verified in 1248/1248 split
primes via Gauss--Cox $\mathrm{Cl}(K)/\mathrm{Cl}(K)^2$ bijection,
cf.\ \cite{OPNewtonDicksonSigned2026}) effects the
$K(\sqrt{d_{\mathrm{genus}}}) \to K \to \Q$ descent that pins each
$q^{(j)}(D)$ in $\Q^\times$.

\medskip
The four-stage programme is:

\begin{enumerate}
\item \emph{Stage 0 (1 week)}: Kings--Sprang Galois-equivariance
check at the smallest $h_K = 2$ anchor $D = -15$. Verify
numerically (PARI/GP, $50$-digit) that the algebraic-period quotient
of $L(\tau \circ \chi, 0)$ lands in
$K(\sqrt{d_{\mathrm{genus}}}, E_\chi)$ as predicted, before any
Q-descent is invoked.

\item \emph{Stage 1 (2 weeks)}: Brunault--Chida $\Q \to K$ descent
for all six invariants. Each $q^{(j)}(D)$ corresponds to a
Rankin--Selberg integral on $Y_0(|D|)(\C)$; the boundary regulator
class lives in $H^1_M(\Spec \Q, M(f_w \otimes f_{w'})(s_w + s_{w'} - w))$,
and the descent to $\Spec K$ is provided by the trace map of
\cite[\S 3]{BrunaultChida2015} applied weight-by-weight.

\item \emph{Stage 2 (2 weeks, principal risk)}: G4 orbit-cancellation
rigorous. The Newton--Dickson signed identity at $h_K = 2$ is
empirically $664/664$ via genus characters (and $1248/1248$ counting
$h_K \in \{1, 2, 4\}$); promoting this to a theorem requires the
explicit Atkin--Lehner per-orbit twist character analysis and a
proof that the cross-Galois cancellation of irrational components
$(\sqrt{d_{\mathrm{genus}}})$ across the two orbits is forced by
the genus-character mechanism.

\item \emph{Stage 3 (1--2 weeks)}: Theorem statement and proof
write-up, $8$--$12$pp arXiv companion to the present paper.
\end{enumerate}

All four ingredients (Kings--Sprang Annals 2025, Kufner 2024,
Brunault--Chida 2015, G4 = OPNewtonDicksonSigned 2026) are in
published / preprint form; no new conjecture is introduced. The
$5$--$7$ week estimate runs in parallel with the Stage 2--5 explicit
$\mathcal{F}$ construction above, and either route alone produces
theorem-level Q-rationality.

\subsection{Categorical generalisations}\label{ss:open-categorical}

A natural generalisation is to drop the constraint $\exp\Cl(K) = 2$
and admit $\exp\Cl(K) | (w - 1)$ for $w \ge 7$, with the Galois group
$(\Z/(w-1))^{r}$. The Newton--Dickson chain (Theorem
\ref{thm:newton-dickson}) extends, but the cross-Galois Q-rationality
of $L$-value polynomials at $w \ge 7$ requires the
Galois-Norm-Multiplicative-Invariant Theorem of
\cite{Paper_P4_2026-05-17}, applied weight-by-weight. Articulating
the full $\mathcal{F}$ on the resulting tower of CM motives is left
to future work.

\subsection{Stark-Heegner structural bound on the empirical anchor
set}\label{ss:heegner-bound}

We record an important structural observation: the empirical pile
underlying Theorems \ref{thm:newton-dickson} and
\ref{thm:beilinson-qD} is intrinsically \emph{finite}, with an
explicit and \emph{closed} list of anchors at class number one.

By the Stark--Heegner theorem (Stark 1967, Heegner 1952), the
imaginary quadratic discriminants $D < 0$ with class number
$h_K = 1$ are exactly the nine values
\[
D \in \{-3,\,-4,\,-7,\,-8,\,-11,\,-19,\,-43,\,-67,\,-163\}.
\]
The Newton--Dickson chain at $h_K = 1$ has $\Z$-coefficients only
in this finite list; for $h_K \ge 2$ the chain takes values in the
CM-unit ring $\Z[\omega]$ of the ring of integers $\mathcal{O}_K$,
and the coefficients descend to $\Z$ only componentwise across the
ideal classes (Saito Galois-descent).

The two cases $D = -3$ and $D = -4$ are special: $\mathcal{O}_K$
has roots of unity of order $6$ and $4$ respectively, and the
associated Hecke Gr\"ossencharacter has trivial character only at
specific weights, breaking the elementary chain. Excluding these
two cases, we have at most \emph{seven} usable Heegner anchors:
\[
\mathcal{H}_{\mathrm{usable}} = \{-7,\,-8,\,-11,\,-19,\,-43,\,-67,\,-163\}.
\]
Combined with the eight $h_K = 2$ anchors of
Theorem \ref{thm:beilinson-qD}, the total empirical baseline for
the present construction is \emph{fifteen} CM discriminants, all
of which we use.

This is a strong structural constraint and not a limitation of
the present technique. The Conjecture $\mathcal{F}$ functor must,
in principle, extend to all $h_K \ge 2$ via Galois descent
componentwise, and to $h_K \to \infty$ via an analytic-continuation
or Iwasawa-tower argument. The fact that the empirical
single-anchor confirmation of Conjecture E-1 motivic-weight
$\leftrightarrow$ glueball-spin (Theorem \ref{thm:su2-ratio}) can
only ever be tested at this fixed list of fifteen anchors is a
\emph{feature}, not a bug: it makes the entire empirical
programme falsifiable in finite time. Either the seven Heegner
$h_K = 1$ anchors satisfy the predicted relations, or they do
not; no infinite-statistics escape is available.

The present paper confirms the seven Heegner anchors at the
Newton--Dickson chain level (Theorem \ref{thm:newton-dickson},
57/57 PARI verifications) and the eight $h_K = 2$ anchors at the
Beilinson $q(D)$ level (Theorem \ref{thm:beilinson-qD}, 8/8 EXACT
to 50 decimal digits). Extension to $h_K \ge 4$ and beyond is
deferred to companion notes.

\subsection{Empirical Faltings-type bound on the 2-torsion
rank}\label{ss:faltings-bound}

For completeness we record an empirical observation relevant to the
extension programme. By the Gauss genus-theory formula
$\mathrm{rk}_2(\mathrm{Cl}(K)) = t - 1$ (where $t$ is the number of
distinct prime divisors of the fundamental discriminant $D$, including
$2$ when $D \equiv 0 \pmod 4$), the 2-torsion rank of $\mathrm{Cl}(K)$
admits a closed-form upper bound from the prime factorisation of $D$
alone. The smallest value of $|D|$ admitting $t = 6$ distinct prime
divisors is
\[
|D| \ge 2^2 \cdot 3 \cdot 5 \cdot 7 \cdot 11 \cdot 13 = 60{,}060,
\]
predicting that no fundamental discriminant with $\mathrm{rk}_2 \ge 5$
exists below this threshold.

A direct PARI/GP sweep of all fundamental discriminants
$D \in [-150{,}000, -3]$ (executed at $60$-digit precision via
\texttt{quadclassunit}) confirms this bound: the total counts are
\begin{center}
\begin{tabular}{l|ccccccc}
$\mathrm{rk}_2(\mathrm{Cl}(K))$ & 0 & 1 & 2 & 3 & 4 & 5 & $\ge 6$\\
\midrule
\#~fundamental disc.\ $|D| \le 150{,}000$ & 5163 & 12930 & 11110 & 3979 & 437 & 6 & \textbf{0}\\
\end{tabular}
\end{center}
The smallest fundamental discriminant with $\mathrm{rk}_2 = 4$ is
$D = -5460 = -2^2 \cdot 3 \cdot 5 \cdot 7 \cdot 13$ (consistent with
the Hasse--Schiitt--Hodge corpus prediction $r(D) = 2^{\mathrm{rk}_2}
= 16$). The smallest fundamental discriminant with
$\mathrm{rk}_2 = 5$ is $D = -87{,}780 = -2^2 \cdot 3 \cdot 5 \cdot 7
\cdot 11 \cdot 19$ (class group structure $[6, 2, 2, 2, 2]$,
$h_K = 96$); five more $\mathrm{rk}_2 = 5$ anchors appear in the
range $|D| \le 150{,}000$ (at $D \in \{-92820, -106260, -120120,
-143220, -145860\}$, each with class-group structure either
$[6, 2, 2, 2, 2]$ or $[8, 2, 2, 2, 2]$). No $\mathrm{rk}_2 \ge 6$
anchor appears in this range; the predicted minimum is
$|D| = 3 \cdot 5 \cdot 7 \cdot 11 \cdot 13 \cdot 17 = 255{,}255$.

Consequently, the Conjecture $\mathcal{F}$ extension programme can
plan in two regimes: a small-$|D|$ regime ($|D| \le 20{,}000$,
characterised by $\mathrm{rk}_2 \le 4$, twenty-one $\mathrm{rk}_2=4$
anchors enumerable on a laptop), and a large-$|D|$ regime requiring
genuine analytic-continuation arguments. The numerics support, but
do not by themselves prove, the analogous Faltings-type bound on
the height of the moduli space of CM K3 surfaces underlying the
arithmetic side of $\mathcal{F}$.

%==================================================================
\section{Conclusion}\label{sec:conclusion}
%==================================================================

This paper articulates a bridge functor $\mathcal{F}$ between the
modern S-matrix bootstrap programme \cite{PaulosI, PaulosII, PaulosIII,
CordovaHeKruczenskiVieira, EliasMiroFluxTube, HeKruczenski23,
HeKruczenski24, KazakovZheng2024}
and the Rankin-Eisenstein motivic $K$-theory of CM modular forms
\cite{Beilinson1985, Schutt2008, BrunaultChida2015, KingsSprang2025},
indexed by an imaginary quadratic field $K = \Q(\sqrt{D})$. A
companion perturbative motivic structure, realised on amplitudes via
the cosmic Galois construction \cite{Brown2016CosmicGalois,
BrownDupont2019, CaronHuotDixon2020}, parallels our non-perturbative
CM-anchored functor (Remark \ref{rk:perturbative-motivic}).

We establish three results unconditionally:
(i) Theorem \ref{thm:beilinson-qD}, the cross-Galois $q(D)$ ratio at all
eight $h_K = 2$ anchors $|D| \le 91$, with denominator $\mid |D|^2$;
(ii) Theorem \ref{thm:newton-dickson}, the Newton--Dickson polynomial
chain for $h_K = 1$ Heegner anchors at all odd weights;
(iii) Theorem \ref{thm:su2-ratio}, the cross-weight motivic ratio
$m_{2^{++}}^2 / m_{0^{++}}^2 = 2$ at the SU(2) lattice anchor
\cite{Athenodorou2021} to $0.07\%$.

These three results constitute the empirical anchors of the
bridge functor, while the construction of $\mathcal{F}$ at the
morphism level is left as a multi-week open programme.

The five mass-gap walls (Section \ref{sec:obstructions}) are all
evaded by the bootstrap-side of $\mathcal{F}$, providing a route
to gauge-theory predictions that does not depend on the
classical constructive frameworks of Balaban, Hairer, Kapustin--Witten
or Polchinski.

\subsection*{Acknowledgements}
We are grateful to the participants of internal discussions during
2026 around the ECI v15 framework and the OP-BEILINSON-Q-D
investigation. Computations were performed in PARI/GP 2.15.4. All
arXiv references in this paper were live-verified via the public
arXiv API on the date of submission.

%==================================================================
\bibliographystyle{amsalpha}
\begin{thebibliography}{99}

\bibitem{Athenodorou2021} A.~Athenodorou, M.~Teper, ``SU(N) gauge
theories in 3+1 dimensions: glueball spectrum, string tensions and
topology,'' \emph{Journal of High Energy Physics} 12 (2021), 082;
arXiv:2106.00364.

\bibitem{Balaban1987} T.~Balaban, ``Renormalization group approach to
lattice gauge field theories,'' \emph{Comm. Math. Phys.} 109 (1987),
no.~2, 249--301.

\bibitem{Beilinson1985} A.~Beilinson, ``Higher regulators and values
of $L$-functions,'' \emph{J.\ Soviet Math.} 30 (1985), 2036--2070.

\bibitem{Bloch1977} S.~Bloch, ``Higher regulators, algebraic
$K$-theory, and zeta functions of elliptic curves,'' Lecture notes,
University of California, Irvine, 1977; published \emph{CRM Monograph
Series} 11, AMS, 2000.

\bibitem{Brown2012} F.~Brown, ``On the decomposition of motivic
multiple zeta values,'' arXiv:1102.1310 (2011); published in
\emph{Adv.\ Stud.\ Pure Math.} 63 (2012), 31--58.

\bibitem{BrunaultChida2015} F.~Brunault, M.~Chida, ``Regulators for
Rankin-Selberg products of modular forms,'' \emph{Annals of
Mathematics Quebec} 40 (2016), 221--249; arXiv:1503.04626.

\bibitem{ChowlaSelberg1949} S.~Chowla, A.~Selberg, ``On Epstein's
zeta function,'' \emph{Proc.\ Nat.\ Acad.\ Sci.\ USA} 35 (1949),
371--374.

\bibitem{CordovaHeKruczenskiVieira} L.~C\'ordova, Y.~He,
M.~Kruczenski, P.~Vieira, ``The O(N) S-matrix Monolith,''
\emph{JHEP} 04 (2020), 142; arXiv:1909.06495.

\bibitem{Cox1989} D.~A.~Cox, \emph{Primes of the form $x^2 + ny^2$},
Wiley-Interscience, 1989.

\bibitem{Damerell1971} R.~M.~Damerell, ``$L$-functions of elliptic
curves with complex multiplication, I,'' \emph{Acta Arithmetica}
17 (1971), 287--301.

\bibitem{EliasMiroFluxTube} J.~Elias-Mir\'o, A.~L.~Guerrieri,
A.~Hebbar, J.~Penedones, P.~Vieira, ``Flux Tube S-matrix Bootstrap,''
\emph{Phys.\ Rev.\ Lett.} 123 (2019), 221602; arXiv:1906.08098.

\bibitem{Gauss1801} C.~F.~Gauss, \emph{Disquisitiones Arithmeticae},
Leipzig, 1801; \S\S 225--237 on genus theory.

\bibitem{Goncharov2001} A.~B.~Goncharov, ``Multiple polylogarithms
and mixed Tate motives,'' arXiv:math/0103059 (2001).

\bibitem{GuerrieriSever} A.~Guerrieri, A.~Sever, ``Rigorous bounds on
the Analytic S-matrix,'' \emph{Phys.\ Rev.\ Lett.} 127 (2021),
251601; arXiv:2106.10257.

\bibitem{Hairer2014} M.~Hairer, ``A theory of regularity
structures,'' \emph{Invent.\ Math.} 198 (2014), 269--504.

\bibitem{HeKruczenski23} Y.~He, M.~Kruczenski, ``Bootstrapping
gauge theories,'' \emph{Phys.\ Rev.\ Lett.} 133 (2024), 191601;
\emph{Phys.\ Rev.\ D} 110 (2024), 096001; arXiv:2309.12402.

\bibitem{HeKruczenski24} Y.~He, M.~Kruczenski, ``Gauge Theory
Bootstrap: Pion amplitudes and low energy parameters,''
arXiv:2403.10772 (2024).

\bibitem{Hecke1937} E.~Hecke, ``\"Uber Modulfunktionen und die
Dirichletschen Reihen mit Eulerscher Produktentwicklung,''
\emph{Math.\ Annalen} 114 (1937), 316--351.

\bibitem{HolleyStroock1989} R.~Holley, D.~Stroock, ``Logarithmic
Sobolev inequalities and stochastic Ising models,''
\emph{J.\ Statist.\ Phys.} 46 (1987), 1159--1194.

\bibitem{KapustinWitten2007} A.~Kapustin, E.~Witten, ``Electric-magnetic
duality and the geometric Langlands program,''
\emph{Commun.\ Number Theory Phys.} 1 (2007), 1--236;
arXiv:hep-th/0604151.

\bibitem{KazakovZheng2024} V.~Kazakov, Z.~Zheng, ``Bootstrap for
finite $N$ lattice Yang--Mills theory,'' arXiv:2404.16925 (2024).

\bibitem{Brown2016CosmicGalois} F.~Brown, ``Feynman amplitudes,
coaction principle, and cosmic Galois group,'' \emph{Communications
in Number Theory and Physics} 11 (2017), 453--556; arXiv:1512.06409.

\bibitem{BrownDupont2019} F.~Brown, C.~Dupont, ``Single-valued
integration and superstring amplitudes in genus zero,''
\emph{Communications in Mathematical Physics} 382 (2021), 815--874;
arXiv:1910.01107.

\bibitem{CaronHuotDixon2020} S.~Caron-Huot, L.~J.~Dixon,
J.~M.~Drummond, F.~Dulat, J.~Foster, \"O.~G\"urdo\u{g}an, M.~von
Hippel, A.~J.~McLeod, G.~Papathanasiou, ``The Steinmann cluster
bootstrap for $\mathcal{N}=4$ super Yang--Mills amplitudes,''
\emph{PoS} CORFU2019 (2020), 003; arXiv:2005.06735.

\bibitem{KingsSprang2025} G.~Kings, J.~Sprang, ``Eisenstein-Kronecker
classes, integrality of critical values of Hecke $L$-functions
and $p$-adic interpolation,'' \emph{Annals of Mathematics} 202
(2025), 1--109; arXiv:1912.03657.

\bibitem{KingsSprangSurvey2025} G.~Kings, J.~Sprang, ``On the
algebraicity and integrality of critical $L$-values of Hecke
characters of imaginary quadratic fields,'' survey, 8pp;
arXiv:2511.05198.

\bibitem{Kufner2024} V.~Kufner, ``On Deligne's conjecture for
algebraic Hecke characters of totally imaginary number fields,''
preprint 2024, 49pp; arXiv:2406.06148.

\bibitem{OPNewtonDicksonSigned2026} K.~R\'emondi\`ere, ``OP-NEWTON-DICKSON-SIGNED
--- Per-orbit twist character resolution via Gauss--Cox genus
mechanism (1248/1248),'' internal note, 2026-05-17.

\bibitem{KoteckyPreiss1986} R.~Koteck\'y, D.~Preiss, ``Cluster
expansion for abstract polymer models,'' \emph{Comm.\ Math.\ Phys.}
103 (1986), 491--498.

\bibitem{OPBeilinsonQD} K.~R\'emondi\`ere, ``OP-BEILINSON-Q-D ---
Beilinson regulator route to a $q(D)$ closed form,'' internal note,
2026-05-17.

\bibitem{OPNewtonDickson2026} K.~R\'emondi\`ere, ``OP-NEWTON-DICKSON-PROOF
--- Rigorous proof of the Newton--Dickson chain for $\Q$-rat CM
newforms,'' internal note, 2026-05-17.

\bibitem{Paper_P4_2026-05-17} K.~R\'emondi\`ere, ``Galois-Norm-Multiplicative
Invariants of CM $L$-values and the Yang-Mills Period Problem,''
Paper P4, 2026-05-17.

\bibitem{PaulosI} M.~F.~Paulos, J.~Penedones, J.~Toledo,
B.~C.~van~Rees, P.~Vieira, ``The S-matrix Bootstrap I: QFT in AdS,''
\emph{JHEP} 11 (2017), 133; arXiv:1607.06109.

\bibitem{PaulosII} M.~F.~Paulos, J.~Penedones, J.~Toledo,
B.~C.~van~Rees, P.~Vieira, ``The S-matrix Bootstrap II: Two
Dimensional Amplitudes,'' \emph{JHEP} 11 (2017), 143;
arXiv:1607.06110.

\bibitem{PaulosIII} M.~F.~Paulos, J.~Penedones, J.~Toledo,
B.~C.~van~Rees, P.~Vieira, ``The S-matrix Bootstrap III: Higher
Dimensional Amplitudes,'' \emph{JHEP} 12 (2019), 040;
arXiv:1708.06765.

\bibitem{PhaseE1} K.~R\'emondi\`ere, ``Phase E --- Extension ECI
pr\'edictions au-del\`a $m_{0^{++}}$,'' internal note 2026-05-15
(memory snapshot \texttt{project\_phase\_e\_eci\_extension\_motivic\_weight}).

\bibitem{Phase8v12} K.~R\'emondi\`ere, ``Phase 8 ECI v15 \'etat net,''
internal note 2026-05-11
(memory snapshot \texttt{project\_phase8\_morn39\_dayend\_v12}).

\bibitem{Polchinski1996} J.~Polchinski, \emph{String Theory},
Cambridge University Press, 1998; vol.~I, \S 6 on flux sectors and
$H^2$.

\bibitem{Ribet1977} K.~A.~Ribet, ``Galois representations attached
to eigenforms with Nebentypus,'' in \emph{Modular Functions of One
Variable V}, LNM 601, Springer, 1977, pp.~17--51.

\bibitem{Scholl1990} A.~J.~Scholl, ``Motives for modular forms,''
\emph{Invent.\ Math.} 100 (1990), 419--430.

\bibitem{Schutt2008} M.~Sch\"utt, ``CM newforms with rational
coefficients,'' \emph{Ramanujan J.} 19 (2009), 187--205;
arXiv:math/0511228.

\bibitem{Shimura1971} G.~Shimura, \emph{Introduction to the arithmetic
theory of automorphic functions}, Iwanami Shoten / Princeton
University Press, 1971.

\bibitem{Shimura1976} G.~Shimura, ``The special values of the zeta
functions associated with cusp forms,'' \emph{Comm.\ Pure Appl.\ Math.}
29 (1976), 783--804.

\bibitem{W3W4PhaseE} K.~R\'emondi\`ere, ``W3+W4 Phase E EXTENSION
FALSIFIED 2026-05-15,'' memory snapshot
\texttt{project\_W3W4\_phase\_e\_extension\_falsified\_2026-05-15}.

\bibitem{BorgaCao2025} P.~Borga, S.~Cao, S.~Shogren, S.~Knaak,
``Surface sums for lattice Yang-Mills in the large-$N$ limit,''
arXiv:2411.11676 (2025).

\bibitem{CEJV2502} E.~Costa, A.~Elsenhans, J.~Jahnel, J.~Voight,
``Kuga-Satake constructions for sextic CM fields,''
arXiv:2502.15052 (2025).

\bibitem{GuerrieriHomrichFluxTube2025} A.~L.~Guerrieri, A.~Homrich,
P.~Vieira, ``Multiparticle Flux Tube S-matrix Bootstrap,''
\emph{Phys.\ Rev.\ Lett.} 134 (2025), 041601; arXiv:2404.10812.

\bibitem{GuerrieriHebbarGlueball2025} A.~L.~Guerrieri, A.~Hebbar,
B.~C.~van Rees, ``Constraining Glueball Couplings,''
\emph{Phys.\ Rev.\ D} 112 (2025), 094023; arXiv:2312.00127.

\end{thebibliography}

\end{document}
